\subsubsection{谱兼容外置、半直积修复与有限/无穷二影刚性}\label{subsubsec:xi-spectrum-compatible-semidirect-two-shadow-rigidity}
本段在 Lee--Yang 四层覆盖的既有单值化框架上，将“可交换外置障碍—非交换修复—线性椭圆双影刚性”组织为同一条可检验链条。

\begin{theorem}[可交换外置寄存器对非交换单值群的不可完备性]\label{thm:xi-terminal-abelian-register-obstruction-nonabelian-monodromy}
设 \(U\) 为道路连通拓扑空间，\(\pi:X\to U\) 为连通有限分支覆盖，分支集记为 \(B\subset U\)，并取基点 \(u_0\in U\setminus B\)。
令
\[
\rho:\pi_1(U\setminus B,u_0)\longrightarrow \mathfrak S_n
\]
为单值表示，且假设 \(\rho(\pi_1(U\setminus B,u_0))\) 非交换。

称三元组 \((A,\iota,D)\) 为该覆盖的一组可交换外置寄存器数据，若：
\begin{enumerate}
  \item \(A\) 为交换群；
  \item \(\iota:\pi_1(U\setminus B,u_0)\to A\) 为群同态；
  \item \(D:\{1,\dots,n\}\times A\to\{1,\dots,n\}\) 满足对任意 \(\gamma\in\pi_1(U\setminus B,u_0)\) 与 \(i\in\{1,\dots,n\}\),
  \[
  \rho(\gamma)(i)=D\bigl(i,\iota(\gamma)\bigr).
  \]
\end{enumerate}
则此类三元组不存在。
\end{theorem}
\begin{proof}
由 \(\rho(\pi_1(U\setminus B,u_0))\) 非交换，可取 \(\alpha,\beta\in\pi_1(U\setminus B,u_0)\) 使
\[
\rho([\alpha,\beta])\neq \mathrm{id},
\qquad
[\alpha,\beta]:=\alpha\beta\alpha^{-1}\beta^{-1}.
\]
由 \(A\) 交换且 \(\iota\) 为同态，
\[
\iota([\alpha,\beta])
=\iota(\alpha)\iota(\beta)\iota(\alpha)^{-1}\iota(\beta)^{-1}
=e_A
=\iota(1).
\]
故对任意 \(i\),
\[
D\bigl(i,\iota([\alpha,\beta])\bigr)=D(i,e_A)=D\bigl(i,\iota(1)\bigr).
\]
左侧按定义等于 \(\rho([\alpha,\beta])(i)\)，右侧等于 \(\rho(1)(i)=i\)，从而 \(\rho([\alpha,\beta])(i)=i\) 对所有 \(i\) 成立，矛盾于 \(\rho([\alpha,\beta])\neq\mathrm{id}\)。
证毕。
\end{proof}

\paragraph{window-6 三轴中心约束、Pati--Salam 连通实现与标准模型包络的中心 2-挠兼容性}
在 window-6 的 boundary 读出中，三轴独立 sheet-flip 给出一个抽象中心寄存器
\[
\mathcal C_{\mathrm{bdry}}\cong (\ZZ/2\ZZ)^3.
\]
本段讨论其与紧连通 Lie 包络在中心 2-挠层面的可实现条件。

\begin{theorem}[标准模型连通包络对三轴中心寄存器的不可实现性]\label{thm:xi-window6-sm-envelope-center-2torsion-obstruction}
设 \(G\) 为紧连通 Lie 群，满足
\[
\Lie(G)\cong \mathfrak u(1)\oplus\mathfrak{su}(2)\oplus\mathfrak{su}(3).
\]
则
\[
\dim_{\FF_2} Z(G)[2]\le 2.
\]
从而不存在单同态
\[
\mathcal C_{\mathrm{bdry}}\cong (\ZZ/2\ZZ)^3\hookrightarrow Z(G)[2].
\]
\end{theorem}
\begin{proof}
由紧连通群分类，可取
\[
G\cong \widetilde G/\Gamma,\qquad
\widetilde G:=U(1)\times SU(2)\times SU(3),
\]
其中 \(\Gamma\subset Z(\widetilde G)\) 为有限中心子群。记
\[
A:=Z(\widetilde G)\cong U(1)\times \ZZ_2\times \ZZ_3,
\qquad
Z(G)\cong A/\Gamma.
\]
设 \(\pi:A\to \ZZ_2\times \ZZ_3\) 为自然投影，则有短正合列
\[
1\to U(1)\to A\to \ZZ_2\times \ZZ_3\to 1.
\]
对 \(\Gamma\) 取商后得到
\[
1\to T\to A/\Gamma\to Q\to 1,
\]
其中 \(T\cong U(1)\)（一维紧环面），\(Q\) 为 \((\ZZ_2\times \ZZ_3)\) 的商群。

记 \(E:=A/\Gamma\)。考虑 \(E[2]\to Q[2]\)。对任意 \(q\in Q[2]\)，其纤维中两点之差落在 \(T[2]\)。
故每个纤维大小至多 \(|T[2]|=2\)，并且
\[
|Q[2]|\le |(\ZZ_2\times \ZZ_3)[2]|=2.
\]
从而
\[
|E[2]|\le 2\cdot 2=4,\qquad
\dim_{\FF_2}E[2]\le 2.
\]
即 \(\dim_{\FF_2} Z(G)[2]\le 2\)。因此 \((\ZZ/2\ZZ)^3\) 不能嵌入 \(Z(G)[2]\)。证毕。
\end{proof}

\begin{theorem}[三轴 canonical sheet-flip 保真下降下的 Pati--Salam 连通实现唯一性]\label{thm:xi-window6-ps-faithful-canonical-descent-unique}
设
\[
\widetilde G_{\mathrm{PS}}:=SU(4)\times SU(2)_L\times SU(2)_R,\qquad
G:=\widetilde G_{\mathrm{PS}}/H,
\]
其中 \(H\subset Z(\widetilde G_{\mathrm{PS}})\) 为有限中心子群。记 canonical involution 子群
\[
\mathcal C_{\mathrm{can}}
:=Z(SU(4))[2]\times Z(SU(2)_L)\times Z(SU(2)_R)
\cong (\ZZ/2\ZZ)^3.
\]
若商映射 \(q:\widetilde G_{\mathrm{PS}}\to G\) 在 \(\mathcal C_{\mathrm{can}}\) 上限制为单同态
\[
q|_{\mathcal C_{\mathrm{can}}}:\mathcal C_{\mathrm{can}}\hookrightarrow Z(G)[2],
\]
则必有 \(H=\{e\}\)，因而
\[
G\cong SU(4)\times SU(2)_L\times SU(2)_R.
\]
\end{theorem}
\begin{proof}
有
\[
Z(\widetilde G_{\mathrm{PS}})
\cong \ZZ_4\times \ZZ_2\times \ZZ_2.
\]
取任意非平凡 \(H\subset Z(\widetilde G_{\mathrm{PS}})\)。任取 \(h\in H\setminus\{e\}\)。
若 \(\ord(h)=2\)，则 \(h\in \mathcal C_{\mathrm{can}}\)；
若 \(\ord(h)=4\)，则 \(h^2\neq e\) 且 \(h^2\in \mathcal C_{\mathrm{can}}\)。
故总有
\[
H\cap \mathcal C_{\mathrm{can}}\neq \{e\}.
\]
而 \(q\) 在 \(H\) 上恒等于单位元，故 \(q|_{\mathcal C_{\mathrm{can}}}\) 的核至少含
\(H\cap \mathcal C_{\mathrm{can}}\) 的非平凡元，不能单射，矛盾。
故 \(H\) 只能为平凡子群。证毕。
\end{proof}

\begin{theorem}[Pati--Salam 到标准模型包络的中心 2-挠压缩不可避免]\label{thm:xi-window6-ps-to-sm-center2-kernel}
设
\[
G_{\mathrm{PS}}:=SU(4)\times SU(2)_L\times SU(2)_R,
\]
且 \(G_{\mathrm{SM}}\) 为任意紧连通 Lie 群，满足
\(
\Lie(G_{\mathrm{SM}})\cong \mathfrak u(1)\oplus\mathfrak{su}(2)\oplus\mathfrak{su}(3)
\)。
若连续群同态
\(
f:G_{\mathrm{PS}}\to G_{\mathrm{SM}}
\)
满足中心兼容条件
\[
f\bigl(\mathcal C_{\mathrm{can}}\bigr)\subseteq Z(G_{\mathrm{SM}})[2],
\]
则限制映射
\[
f_*:\mathcal C_{\mathrm{can}}\longrightarrow Z(G_{\mathrm{SM}})[2]
\]
必有非平凡核。
\end{theorem}
\begin{proof}
由定理 \ref{thm:xi-window6-sm-envelope-center-2torsion-obstruction}，
\(
\dim_{\FF_2} Z(G_{\mathrm{SM}})[2]\le 2
\)。
而
\(
\mathcal C_{\mathrm{can}}\cong (\ZZ/2\ZZ)^3
\)
的 \(\FF_2\)-维数为 \(3\)。
故线性映射 \(f_*\) 的秩至多 \(2\)，其核维数至少 \(1\)，即核非平凡。证毕。
\end{proof}

\begin{proposition}[sheet parity 在 \(m=6\) 的严格锚点性]\label{prop:xi-window6-sheet-parity-anchor-m6}
\leanverified{paper\_xi\_window6\_sheet\_parity\_anchor\_m6}
在 dyadic boundary uplift 口径下，取
\[
\Delta_m^{\mathrm{bdry}}
:=\{n-V_m(w):\ n\in(\Fold_m^{\mathrm{bin}})^{-1}(w)\}
\]
（对 boundary \(w\) 不依赖于具体词）。
核验给出
\[
\Delta_6^{\mathrm{bdry}}=\{0,34\},\qquad
\Delta_7^{\mathrm{bdry}}=\{0,55,89\},\qquad
\Delta_8^{\mathrm{bdry}}=\{0,89,144\}.
\]
在“parity 由 uplift 数据无需额外选择即可恢复”的准则下，该结构仅在 \(m=6\) 成立。
\end{proposition}
\begin{proof}
若 parity 可由 \(\Delta_m^{\mathrm{bdry}}\) 唯一恢复，则每个 boundary 纤维必须存在唯一二元配对规则；
这要求非零 uplift 增量唯一。
当 \(m=6\) 时，\(\Delta_6^{\mathrm{bdry}}\) 仅含一个非零元 \(34\)，且 boundary 预像层数为 \(2\)，配对唯一。
当 \(m=7,8\) 时，\(\Delta_m^{\mathrm{bdry}}\) 各含两个不同非零元且预像层数为 \(3\)；
任意二元化都需额外选择（例如忽略某一层），不再由数据唯一决定。证毕。
\end{proof}

\begin{theorem}[纤维内自由 involution 的奇数障碍与最小覆盖度]\label{thm:xi-bdry-free-involution-odd-fiber-obstruction-min-cover}
\leanverified{paper\_xi\_bdry\_free\_involution\_odd\_fiber\_obstruction\_min\_cover}
设 \(f:\widetilde X\to X\) 为有限集合之间的满射。
若存在自由 involution \(\sigma:\widetilde X\to\widetilde X\) 满足
\[
\sigma^2=\mathrm{id},\qquad
\mathrm{Fix}(\sigma)=\varnothing,\qquad
f\circ\sigma=f,
\]
则对任意 \(x\in X\) 必有
\[
\bigl|f^{-1}(x)\bigr|\equiv 0\pmod 2.
\]
因此若存在 \(x\in X\) 使 \(|f^{-1}(x)|\) 为奇数，则不存在满足 \(f\circ\sigma=f\) 的自由 involution。
\par
进一步，若希望通过一个覆盖 \(\pi:Y\to\widetilde X\) 在 \(Y\) 上实现同类的纤维内自由 involution（即 \((f\circ\pi)\circ\widetilde\sigma=f\circ\pi\) 且 \(\widetilde\sigma\) 自由），
则覆盖度必须为偶数，因而 \(\deg(\pi)\ge 2\)。
\end{theorem}
\begin{proof}
若 \(\sigma\) 自由，则其轨道分解全由二元轨道组成，故 \(|S|\) 为偶数对任何 \(\sigma\)-不变有限子集 \(S\subseteq \widetilde X\) 成立。
取 \(S=f^{-1}(x)\)。由 \(f\circ\sigma=f\) 得 \(S\) 为 \(\sigma\)-不变，故 \(|f^{-1}(x)|\) 必为偶数；若存在奇数纤维则矛盾。
\par
对覆盖情形：若 \(|f^{-1}(x)|\) 为奇数，则 \(|(f\circ\pi)^{-1}(x)|=\deg(\pi)\cdot |f^{-1}(x)|\) 与 \(\deg(\pi)\) 同奇偶。
而在 \(Y\) 上实现纤维内自由 involution 仍要求每条纤维基数为偶数，故 \(\deg(\pi)\) 必为偶数，特别地 \(\deg(\pi)\ge 2\)。证毕。
\end{proof}

\begin{corollary}[boundary 三层 uplift 的二值外置下界]\label{cor:xi-bdry-uplift-three-layer-needs-one-bit}
在 dyadic boundary uplift 口径下，若 \(m\in\{7,8\}\)，则由命题 \ref{prop:xi-window6-sheet-parity-anchor-m6} 的
\(
\Delta_m^{\mathrm{bdry}}=\{0,\cdot,\cdot\}
\)
可知 boundary 预像层数为 \(3\)，从而存在 boundary 词 \(w\) 使
\[
\bigl|(\Fold_m^{\mathrm{bin}})^{-1}(w)\bigr|=3.
\]
因此不存在在 boundary uplift 纤维内闭合的自由 \(\ZZ_2\) involution；
并且任何试图在覆盖替代后实现该自由性的数据结构，其最小覆盖度满足 \(\deg(\pi)\ge 2\)，从而至少需要引入一个二值外置轴以消除奇数纤维的轨道障碍。
\end{corollary}
\begin{proof}
由命题 \ref{prop:xi-window6-sheet-parity-anchor-m6}，当 \(m=7,8\) 时 boundary 预像层数为 \(3\)，故存在边界纤维基数为 \(3\)。
将定理 \ref{thm:xi-bdry-free-involution-odd-fiber-obstruction-min-cover} 应用于 \(f=(\Fold_m^{\mathrm{bin}})|_{\widetilde X_m^{\mathrm{bdry}}}\) 即得结论。证毕。
\end{proof}

\begin{theorem}[全对称不变性强制的 \texorpdfstring{$\ZZ_2$}{Z/2} 分级平凡化]\label{thm:xi-sym-invariant-z2-grading-trivial}\leanverified{paper\_xi\_sym\_invariant\_z2\_grading\_trivial}
设 \(\Delta\) 为有限集合，\(|\Delta|=k\ge 3\)。
将一个 \(\ZZ_2\)-分级理解为选取“奇层子集” \(A\subset\Delta\)（允许 \(A=\varnothing\) 或 \(A=\Delta\) 的平凡分级）。
若要求该分级在全对称群 \(\mathrm{Sym}(\Delta)\cong S_k\) 的自然作用下不变，即对任意 \(\sigma\in S_k\) 都有 \(\sigma(A)=A\)，则该分级必为平凡分级：\(A=\varnothing\) 或 \(A=\Delta\)。
\end{theorem}
\begin{proof}
由于 \(S_k\) 在 \(\Delta\) 上传递，若 \(A\neq\varnothing\)，取 \(x\in A\)。
对任意 \(y\in\Delta\) 存在 \(\sigma\in S_k\) 使 \(\sigma(x)=y\)。
由不变性 \(\sigma(A)=A\) 得 \(y\in A\)，从而 \(A=\Delta\)。
同理，若 \(A\neq\Delta\) 则只能 \(A=\varnothing\)。
故除平凡分级外不存在 \(S_k\)-不变的 \(\ZZ_2\) 分级。证毕。
\end{proof}

\begin{theorem}[左--右交换下的规范中心模与不可下降通道]\label{thm:xi-window6-lr-swap-canonical-center-mod}
在 \(\mathcal C_{\mathrm{can}}\cong(\ZZ/2\ZZ)^3\) 中取基
\[
a:=(-I_4,1,1),\qquad b:=(1,-I_2,1),\qquad c:=(1,1,-I_2).
\]
令 \(\iota\) 为交换 \(SU(2)_L\leftrightarrow SU(2)_R\) 的自同构，则
\[
\iota(a)=a,\qquad \iota(b)=c,\qquad \iota(c)=b.
\]
定义满同态
\[
\pi_{\mathrm{LR}}:\mathcal C_{\mathrm{can}}\twoheadrightarrow \ZZ/2\ZZ,\qquad
\pi_{\mathrm{LR}}(\alpha a+\beta b+\gamma c):=\beta+\gamma.
\]
则
\[
\ker(\pi_{\mathrm{LR}})=\mathcal C_{\mathrm{can}}^{\iota}
\cong (\ZZ/2\ZZ)^2.
\]
因而任何要求“左--右交换在中心层不可见”的下降机制，必然至少压缩一个规范 \(\ZZ_2\) 模通道。
\end{theorem}
\begin{proof}
\(\pi_{\mathrm{LR}}\) 显然为群同态且满。
\[
\pi_{\mathrm{LR}}(\alpha,\beta,\gamma)=0
\iff \beta=\gamma.
\]
另一方面，
\[
\iota(\alpha,\beta,\gamma)=(\alpha,\gamma,\beta),
\]
故
\(
(\alpha,\beta,\gamma)\in\mathcal C_{\mathrm{can}}^{\iota}
\iff \beta=\gamma
\)。
两条件等价，故核恰为不动子群，且其维数为 \(2\)。证毕。
\end{proof}

\begin{conjecture}[dyadic boundary uplift 层数在 \(m=6\) 的唯一退化]\label{conj:xi-window6-bdry-uplift-unique-degeneration}
设 \(\Delta_m^{\mathrm{bdry}}\) 为 boundary uplift 的增量集合。猜想对所有 \(m\ge 7\) 有
\[
\Delta_m^{\mathrm{bdry}}=\{0,F_{m+3},F_{m+4}\},
\]
从而 boundary 覆盖层数恒为 \(3\)。
而 \(m=6\) 为唯一退化分辨率：
\[
\Delta_6^{\mathrm{bdry}}=\{0,F_9\}=\{0,34\},
\]
即唯一仅含单一非零增量的二层情形。
\end{conjecture}

\begin{conjecture}[异常与共振的通道稀疏性：由 boundary sheet parity 子群支配]\label{conj:xi-anomaly-resonance-channel-sparsity-controlled-by-boundary-parity}
在 window-\(6\) 的 boundary 读出中，三轴独立 sheet-flip 诱导规范中心寄存器
\(\mathcal C_{\mathrm{bdry}}\cong(\ZZ/2\ZZ)^3\)。
猜想：对所有足够大的分辨率 \(m\)，所有可见通道中的非平凡绕数/真共振指数（从而异常的可见非零性）只能出现在那些在 \(\mathcal C_{\mathrm{bdry}}\) 上取非平凡限制的特征标通道中。
等价地：若某一可见通道 \(\overline\chi\) 在 \(\mathcal C_{\mathrm{bdry}}\) 上平凡，则其可见散射绕数与相应的 Toeplitz 负平方指数均为零。
\par
该稀疏性一旦成立，将把“通道级异常定位/证书化”的搜索空间压缩到由 \(\mathcal C_{\mathrm{bdry}}\) 诱导的有限层筛选子族，并与定理 \ref{thm:xi-anom-winding-obstruction} 与定理 \ref{thm:xi-visible-anom-iff-radial-wallcrossing} 的可见通道判别口径形成闭合接口。
\end{conjecture}

\begin{remark}[可复现核验脚本]\label{rem:xi-window6-center-sheetflip-ps-sm-audit-script}
本段代数与组合核验由脚本
\texttt{scripts/exp\_xi\_window6\_center\_sheetflip\_ps\_sm\_audit.py}
给出，证书文件为
\texttt{artifacts/export/xi\_window6\_center\_sheetflip\_ps\_sm\_audit.json}。
\end{remark}


\begin{corollary}[Lee--Yang 四层覆盖中的可交换外置障碍]\label{cor:xi-terminal-abelian-register-obstruction-leyang-s4}
在定理 \ref{thm:xi-terminal-zm-leyang-elliptic-structure} 给出的 Lee--Yang 四层覆盖 \(\pi_y:E_{\min}\to\mathbb P^1_y\) 上，
\[
\mathrm{Mon}(\pi_y)\cong S_4
\]
为非交换群，故满足定理 \ref{thm:xi-terminal-abelian-register-obstruction-nonabelian-monodromy} 的不可完备性结论：任何要求路径拼接同态化的可交换外置寄存器都不能实现全局单值恢复。
\end{corollary}

\begin{theorem}[位移半直积 Gödel 寄存器对非交换算法轨迹的忠实外置]\label{thm:xi-terminal-semidir-godel-register-faithful}
\leanverified{paper\_xi\_terminal\_semidir\_godel\_register\_faithful}
设 \(\Sigma\) 为有限字母表（\(|\Sigma|\ge 2\)），\(\Sigma^\ast\) 为自由幺半群（串联为乘法），并取单射编码
\(
c:\Sigma\hookrightarrow \NN_{\ge 1}
\)。
记 \((p_i)_{i\ge 1}\) 为递增素数序列。

对 \(w=a_1\cdots a_n\in\Sigma^\ast\) 定义位序 Gödel 编码
\[
G(w):=\prod_{i=1}^{n}p_i^{\,c(a_i)},
\qquad
|w|=n.
\]
定义素数位移同态
\[
\mathrm{Sh}^{k}\!\left(\prod_{i\ge1}p_i^{e_i}\right):=\prod_{i\ge1}p_{i+k}^{e_i}.
\]
令
\[
\mathcal R:=\NN\times \NN_{>0},
\qquad
(n,A)\odot(m,B):=\bigl(n+m,\ A\cdot \mathrm{Sh}^{n}(B)\bigr).
\]
则映射
\[
J:\Sigma^\ast\to\mathcal R,\qquad
J(w):=(|w|,G(w))
\]
为单射幺半群同态，且 \((|w|,G(w))\) 可唯一恢复 \(w\)。
此外，对任意交换幺半群 \(M\)，不存在单射幺半群同态 \(\Sigma^\ast\hookrightarrow M\)。
\end{theorem}
\begin{proof}
\(\mathrm{Sh}^{n}\circ\mathrm{Sh}^{m}=\mathrm{Sh}^{n+m}\) 且每个 \(\mathrm{Sh}^n\) 为幺半群同态，故 \(\odot\) 结合，单位元为 \((0,1)\)。
对 \(u=a_1\cdots a_n\)、\(v=b_1\cdots b_m\)，
\[
G(uv)
=\left(\prod_{i=1}^{n}p_i^{c(a_i)}\right)\left(\prod_{j=1}^{m}p_{n+j}^{c(b_j)}\right)
=G(u)\,\mathrm{Sh}^{n}(G(v)),
\]
故 \(J(uv)=J(u)\odot J(v)\)。

单射性由唯一分解定理给出：\(J(w)=(n,G(w))\) 已给出长度 \(n\)，且
\(
v_{p_i}(G(w))=c(a_i)
\)
可逐位恢复 \(a_i\)（因 \(c\) 单射）。

若 \(f:\Sigma^\ast\to M\) 为到交换幺半群的同态，取 \(a\neq b\in\Sigma\)，则
\[
f(ab)=f(a)f(b)=f(b)f(a)=f(ba).
\]
若 \(f\) 单射则推出 \(ab=ba\)，与自由幺半群非交换性矛盾。证毕。
\end{proof}

% AUTO-GENERATED by scripts/exp_xi_prime_register_semidirect_two_shadow_audit.py
\begin{equation}\label{eq:xi_prime_register_semidirect_two_shadow_audit}
\begin{aligned}
\rho([\alpha,\beta])&=(2,3,1,4),\qquad \rho([\alpha,\beta])(1)\neq 1,\\
J(uv)&=J(u)\odot J(v)\ \text{(checked for all words with }|u|,|v|\le 4\text{)},\\
\left|O_2(\ZZ)\right|&=8,\qquad \left|SO_2(\ZZ)\right|=4.
\end{aligned}
\end{equation}


\begin{theorem}[Lee--Yang--Gödel 嵌入与素数寄存器的代数可识别性]\label{thm:xi-leyang-godel-prime-register-algebraic-identifiability}
令 \(\mathcal P\) 为素数集。记
\[
\NN^{(\mathcal P)}:=\{r:\mathcal P\to\NN:\ \#\mathrm{supp}(r)<\infty\},
\qquad
\mathrm{supp}(r):=\{p\in\mathcal P:\ r(p)\neq 0\}.
\]
并定义 Gödel 读数
\[
N(r):=\prod_{p\in\mathcal P}p^{r(p)}\in\NN_{\ge 1}.
\]
定义映射
\[
\mathcal L:\NN^{(\mathcal P)}\to \ZZ[z],
\qquad
\mathcal L(r)(z):=\prod_{p\in\mathcal P}\bigl(1+z^p\bigr)^{r(p)}.
\]
则成立：
\begin{enumerate}
  \item \(\mathcal L\) 为交换幺半群同态（以点加 \(r+s\) 与多项式乘法为幺半群运算）：
  \[
  \mathcal L(r+s)=\mathcal L(r)\mathcal L(s).
  \]
  \item \(\mathcal L\) 为单射。
  \item \(\mathcal L(r)\) 的全部零点均落在单位圆 \(S^1\)，且皆为偶阶单位根。
\end{enumerate}
\end{theorem}
\begin{proof}
第一条由指数加法与乘法分配律直接给出。
\par
第三条：任一因子 \(1+z^p\) 的零点满足 \(z^p=-1\)，从而 \(z^{2p}=1\)，故其零点均为偶阶单位根，亦在 \(S^1\) 上；乘积的零点集为各因子零点的并（按重数计），结论成立。
\par
第二条：对奇素数 \(p\) 有分解
\[
z^p+1=(z+1)\Phi_{2p}(z),
\]
且 \(\Phi_{2p}\) 在 \(\QQ\) 上不可约，并且不同 \(p\) 给出不同的不可约因子（可参见 \cite{Serre1973CourseArithmetic} 对 cyclotomic 多项式的基本性质与不可约性论述）。
对 \(p=2\)，有 \(z^2+1=\Phi_4(z)\)。
因此在 \(\mathcal L(r)\) 的不可约分解中，每个 \(\Phi_{2p}\)（\(p\) 为奇素数）出现的指数恰为 \(r(p)\)，而 \(\Phi_4\) 出现的指数为 \(r(2)\)，从而 \(r\) 由 \(\mathcal L(r)\) 唯一确定，故 \(\mathcal L\) 单射。
\end{proof}

\begin{theorem}[Gödel--Lee--Yang 对数振幅的逐素频全息反演]\label{thm:xi-leyang-logamplitude-prime-frequency-inversion}
在定理 \ref{thm:xi-leyang-godel-prime-register-algebraic-identifiability} 的记号下，对任意 \(r\in\NN^{(\mathcal P)}\) 定义单位圆上的对数振幅
\[
\mathcal A_r(\theta):=\log\abs{\mathcal L(r)(e^{i\theta})}\qquad(\theta\in[-\pi,\pi]).
\]
则 \(\mathcal A_r\in L^1([-\pi,\pi])\)。
并且对每个素数 \(p\) 成立严格反演公式
\[
\boxed{\quad
r(p)=\frac1\pi\int_{-\pi}^{\pi}\mathcal A_r(\theta)\cos(p\theta)\,\dd\theta.\quad}
\]
因此 \(\{\mathcal A_r(\theta)\}_{\theta\in[-\pi,\pi]}\) 在测度意义下逐素数唯一确定寄存器 \(r\)。
\end{theorem}
\begin{proof}
由于 \(r\) 仅有限支撑，
\[
\mathcal A_r(\theta)=\sum_{p\in\mathcal P}r(p)\log\abs{1+e^{ip\theta}}.
\]
对任意 \(\rho\in(0,1)\)，在单位圆盘内取主值对数并取实部，得到一致可积展开
\[
\Re\log(1+\rho e^{ix})
=\sum_{k\ge 1}\frac{(-1)^{k+1}}{k}\rho^k\cos(kx).
\]
令 \(\rho\uparrow 1\) 可得在 \(L^2([-\pi,\pi])\)（从而 \(L^1\)）意义下的余弦级数
\[
\log\abs{1+e^{ix}}
=\sum_{k\ge 1}\frac{(-1)^{k+1}}{k}\cos(kx)
\qquad(x\in[-\pi,\pi]\text{ a.e.}).
\]
代入 \(x=p\theta\) 并对有限个 \(p\) 求和（可交换求和次序）得到
\[
\mathcal A_r(\theta)
=\sum_{p}r(p)\sum_{k\ge 1}\frac{(-1)^{k+1}}{k}\cos(kp\theta).
\]
因此 \(\mathcal A_r\) 的第 \(n\) 个余弦系数为
\[
\frac1\pi\int_{-\pi}^{\pi}\mathcal A_r(\theta)\cos(n\theta)\,\dd\theta
=\sum_{p\mid n} r(p)\,\frac{(-1)^{n/p+1}}{n/p}.
\]
当 \(n=p\) 为素数时右端仅含 \(p\) 一项且 \(n/p=1\)，从而该余弦系数等于 \(r(p)\)，即得所述反演公式。证毕。
\end{proof}

\begin{corollary}[整数刚性阈值：\(L^2\) 误差控制下的逐素数精确恢复]\label{cor:xi-leyang-logamplitude-rounding-threshold}
在定理 \ref{thm:xi-leyang-logamplitude-prime-frequency-inversion} 的设定下，令 \(\widetilde{\mathcal A}\in L^2([-\pi,\pi])\) 为任意观测/学习得到的对数振幅近似，并对素数 \(p\) 定义
\[
\widetilde r(p):=\operatorname{Round}\!\left(\frac1\pi\int_{-\pi}^{\pi}\widetilde{\mathcal A}(\theta)\cos(p\theta)\,\dd\theta\right)\in\ZZ,
\]
其中 \(\operatorname{Round}\) 表示四舍五入到最近整数。
若满足误差阈值
\[
\norm{\widetilde{\mathcal A}-\mathcal A_r}_{L^2([-\pi,\pi])}<\frac{\sqrt\pi}{2},
\]
则对所有素数 \(p\) 均有 \(\widetilde r(p)=r(p)\)。
\end{corollary}
\begin{proof}
记
\[
c_p(F):=\frac1\pi\int_{-\pi}^{\pi}F(\theta)\cos(p\theta)\,\dd\theta.
\]
由 Cauchy--Schwarz 不等式并注意 \(\norm{\cos(p\theta)}_{L^2([-\pi,\pi])}=\sqrt\pi\)，得到
\[
\abs{c_p(\widetilde{\mathcal A})-c_p(\mathcal A_r)}
\le \frac1\pi\norm{\widetilde{\mathcal A}-\mathcal A_r}_{L^2}\,\norm{\cos(p\theta)}_{L^2}
=\frac1{\sqrt\pi}\norm{\widetilde{\mathcal A}-\mathcal A_r}_{L^2}
<\frac12.
\]
而由定理 \ref{thm:xi-leyang-logamplitude-prime-frequency-inversion}，\(c_p(\mathcal A_r)=r(p)\in\ZZ\)。
因此 \(c_p(\widetilde{\mathcal A})\) 与整数 \(r(p)\) 的距离严格小于 \(1/2\)，四舍五入必回到 \(r(p)\)。证毕。
\end{proof}

\begin{definition}[素数根均匀测度与零点经验测度]\label{def:xi-leyang-prime-root-measures}
对每个素数 \(p\)，令 \(\eta_p\) 为方程 \(z^p=-1\) 的 \(p\) 个解上的均匀概率测度：
\[
\eta_p:=\frac1p\sum_{k=0}^{p-1}\delta_{\exp\!\left(i\frac{(2k+1)\pi}{p}\right)}
\quad\text{（概率测度，支撑于 \(S^1\)）}.
\]
对任意 \(r\in\NN^{(\mathcal P)}\setminus\{0\}\)，记
\[
D(r):=\deg \mathcal L(r)=\sum_{p\in\mathcal P}r(p)\,p,
\]
并定义其零点归一化经验测度（按重数计）
\[
\mu_r:=\frac1{D(r)}\sum_{\zeta:\ \mathcal L(r)(\zeta)=0}\mathrm{mult}(\zeta)\,\delta_\zeta
\quad\text{（概率测度，支撑于 \(S^1\)）}.
\]
\end{definition}

\begin{lemma}[Fourier 系数的显式消失律]\label{lem:xi-eta-p-fourier}
\leanverified{paper\_xi\_eta\_p\_fourier}
对任意整数 \(n\neq 0\)，有
\[
\widehat{\eta_p}(n):=\int_{S^1}z^n\,\dd\eta_p(z)=
\begin{cases}
(-1)^{n/p},&p\mid n,\\
0,&p\nmid n.
\end{cases}
\]
\end{lemma}
\begin{proof}
\[
\int_{S^1}z^n\,\dd\eta_p(z)
=\frac1p\sum_{k=0}^{p-1}\exp\!\left(i n\frac{(2k+1)\pi}{p}\right)
=e^{i n\pi/p}\cdot \frac1p\sum_{k=0}^{p-1}e^{i2\pi nk/p}.
\]
若 \(p\nmid n\)，则内层为非平凡 \(p\) 次单位根之和，取值 \(0\)；若 \(p\mid n\)，则内层为 \(1\)，从而整体为 \(e^{i\pi(n/p)}=(-1)^{n/p}\)。
\end{proof}

\begin{proposition}[零点测度的素数分解与单频谱解码]\label{prop:xi-leyang-prime-measure-linear-independence-decoding}
对任意 \(r\in\NN^{(\mathcal P)}\setminus\{0\}\)，有凸组合分解
\[
\mu_r=\sum_{p\in\mathcal P}\frac{r(p)\,p}{D(r)}\,\eta_p.
\]
进一步，对任意有限支撑整数族 \((c_p)_{p\in\mathcal P}\subset\ZZ\) 与带符号测度
\(
\sigma:=\sum_pc_p\,\eta_p
\)，若 \(\sigma=0\) 则必有 \(c_p=0\)（对所有 \(p\)）。
并且对任意素数 \(p\)，有解码公式
\[
c_p=-\int_{S^1}z^p\,\dd\sigma(z).
\]
\end{proposition}
\begin{proof}
由定义 \(\mathcal L(r)=\prod_p(1+z^p)^{r(p)}\)，每个因子 \(1+z^p\) 贡献 \(p\) 个单位根零点，且零点重数按指数 \(r(p)\) 线性叠加，因此零点测度满足
\(
\mu_r=\sum_p\frac{r(p)p}{D(r)}\eta_p
\)。
\par
对线性无关性与解码，取 Fourier 矩 \(n=p\)。由引理 \ref{lem:xi-eta-p-fourier}，对 \(q\neq p\) 因 \(q\nmid p\) 有
\(
\int z^p\,\dd\eta_q=0
\)，而对 \(q=p\) 有
\(
\int z^p\,\dd\eta_p=-1
\)。
故
\[
\int_{S^1}z^p\,\dd\sigma(z)=\sum_q c_q\int z^p\,\dd\eta_q(z)=-c_p,
\]
从而若 \(\sigma=0\) 则 \(c_p=0\)（对任意 \(p\)），并得解码公式。
\end{proof}

\begin{theorem}[素数截断零点谱的 Haar 极限与显式差分界]\label{thm:xi-leyang-prime-truncation-haar-discrepancy}
令
\[
\mathcal L_{\le P}(z):=\prod_{p\le P}(1+z^p),
\qquad
D_P:=\deg\mathcal L_{\le P}=\sum_{p\le P}p,
\]
并令 \(\mu_P\) 为 \(\mathcal L_{\le P}\) 的零点归一化经验测度（按重数计）。
则当 \(P\to\infty\) 时，\(\mu_P\) 在弱-\(\ast\) 意义下收敛到单位圆上的 Haar 概率测度 \(m_{S^1}\)。
此外，存在绝对常数 \(C>0\)，使对任意弧段 \(A\subset S^1\) 有
\[
\sup_{A\subset S^1}\abs{\mu_P(A)-m_{S^1}(A)}
\le \frac{C}{\sqrt{D_P}}
=\frac{C}{\sqrt{\sum_{p\le P}p}}.
\]
\end{theorem}
\begin{proof}
由命题 \ref{prop:xi-leyang-prime-measure-linear-independence-decoding} 的分解（取 \(r(p)=\ind_{\{p\le P\}}\)），有
\[
\mu_P=\sum_{p\le P}\frac{p}{D_P}\,\eta_p.
\]
对任意整数 \(n\neq 0\)，其 Fourier 系数为
\[
\widehat{\mu_P}(n)
=\sum_{p\le P}\frac{p}{D_P}\widehat{\eta_p}(n)
=\frac1{D_P}\sum_{\substack{p\le P\\ p\mid n}}p\,(-1)^{n/p},
\]
其中最后一步用引理 \ref{lem:xi-eta-p-fourier}。
对固定 \(n\neq 0\)，分子仅为有限和且与 \(P\) 无关，而 \(D_P\to\infty\)，故 \(\widehat{\mu_P}(n)\to 0\)。
Haar 测度 \(m_{S^1}\) 的所有非零 Fourier 系数均为 \(0\)，因而 \(\mu_P\Rightarrow m_{S^1}\)。
\par
差分界由 Erd\H{o}s--Tur\'an 不等式给出（参见 \cite{KuipersNiederreiter1974}）：存在绝对常数 \(C_0>0\)，对任意 \(H\in\NN\) 有
\[
\sup_A\abs{\mu_P(A)-m_{S^1}(A)}
\le C_0\left(\frac1H+\sum_{k=1}^H\frac{|\widehat{\mu_P}(k)|}{k}\right).
\]
由上式表达与 \(\sum_{p\mid k}p\le k\) 得
\(
|\widehat{\mu_P}(k)|\le k/D_P
\)，从而
\[
\sum_{k=1}^H\frac{|\widehat{\mu_P}(k)|}{k}\le \sum_{k=1}^H\frac1{D_P}=\frac{H}{D_P}.
\]
取 \(H=\lfloor\sqrt{D_P}\rfloor\) 得
\(
\sup_A\abs{\mu_P(A)-m_{S^1}(A)}\le 2C_0/\sqrt{D_P}
\)，并将常数并入 \(C\) 即得。
\end{proof}

\begin{proposition}[Gödel 椭圆边界的推前测度与显式密度]\label{prop:xi-godel-ellipse-pushforward-density}
取 \(a,b>0\)，并定义单位圆到轴对齐椭圆边界的参数化
\[
T_{a,b}:S^1\to \partial E_{a,b},
\qquad
T_{a,b}(e^{it})=(a\cos t,\ b\sin t),
\]
其中 \(\partial E_{a,b}\) 为 \(x^2/a^2+y^2/b^2=1\) 的边界。
令
\(
\nu_{a,b}:=(T_{a,b})_*m_{S^1}
\)
为 Haar 测度的推前测度，则 \(\nu_{a,b}\) 关于椭圆弧长测度 \(ds\) 绝对连续，且 Radon--Nikodym 导数满足
\[
\frac{\dd\nu_{a,b}}{\dd s}\Bigl(T_{a,b}(e^{it})\Bigr)
=\frac1{2\pi}\cdot\frac1{\sqrt{a^2\sin^2 t+b^2\cos^2 t}}.
\]
进一步，若 \(\mu_P\) 为定理 \ref{thm:xi-leyang-prime-truncation-haar-discrepancy} 中的零点谱测度，则
\[
(T_{a,b})_*\mu_P\Rightarrow \nu_{a,b}.
\]
\end{proposition}
\begin{proof}
由弧长元素
\(
ds=\sqrt{(a\sin t)^2+(b\cos t)^2}\,\dd t
\)
可知 \(m_{S^1}\) 在参数 \(t\) 下的密度为 \(\dd t/(2\pi)\)。推前到 \(\partial E_{a,b}\) 后得到关于 \(ds\) 的密度为
\[
\frac{\dd t}{2\pi}\cdot \frac1{\dd s}
=\frac1{2\pi}\cdot\frac1{\sqrt{a^2\sin^2 t+b^2\cos^2 t}}.
\]
弱收敛在连续映射下的推前稳定性给出第二条：由 \(\mu_P\Rightarrow m_{S^1}\) 且 \(T_{a,b}\) 连续，得 \((T_{a,b})_*\mu_P\Rightarrow (T_{a,b})_*m_{S^1}=\nu_{a,b}\)。
\end{proof}

\begin{theorem}[素数寄存器序列的 Haar 光谱收敛判别]\label{thm:xi-leyang-prime-register-haar-criterion}
设 \((r_m)_{m\ge 1}\subset \NN^{(\mathcal P)}\setminus\{0\}\) 为任意序列，并令
\[
\mathcal L_m(z):=\mathcal L(r_m)(z),
\qquad
D_m:=D(r_m)=\sum_{p\in\mathcal P}r_m(p)\,p,
\qquad
\mu_m:=\mu_{r_m}.
\]
则以下三条等价：
\begin{enumerate}
  \item \(\mu_m\Rightarrow m_{S^1}\)。
  \item 对每个固定素数 \(q\)，有
  \[
  \frac{r_m(q)\,q}{D_m}\longrightarrow 0.
  \]
  \item 对每个固定整数 \(n\neq 0\)，有
  \[
  \frac{\sum_{p\mid n}r_m(p)\,p}{D_m}\longrightarrow 0.
  \]
\end{enumerate}
\end{theorem}
\begin{proof}
由定义 \(\mu_m=\sum_p \frac{r_m(p)p}{D_m}\eta_p\) 与引理 \ref{lem:xi-eta-p-fourier}，对任意整数 \(n\neq 0\) 有
\[
\widehat{\mu_m}(n)
=\int z^n\,\dd\mu_m(z)
=\frac1{D_m}\sum_{p\mid n}r_m(p)\,p\,(-1)^{n/p}.
\]
Haar 测度满足 \(\widehat{m_{S^1}}(n)=0\)（\(n\neq 0\)），故 \(\mu_m\Rightarrow m_{S^1}\) 等价于对一切 \(n\neq 0\) 有 \(\widehat{\mu_m}(n)\to 0\)；
而这与第三条等价（符号项不影响绝对值阶）。
第三条显然推出第二条（取 \(n=q\)）。
第二条推出第三条是因为 \(n\) 的素因子集合有限，\(\sum_{p\mid n}r_m(p)p/D_m\) 为有限项之和，逐项趋于 \(0\) 即整体趋于 \(0\)。
\end{proof}

\begin{theorem}[零点普适律与对数振幅保持律的严格分裂]\label{thm:xi-leyang-bulk-boundary-splitting-haar-vs-logamplitude}
设 \((r_N)_{N\ge 1}\subset \NN^{(\mathcal P)}\setminus\{0\}\)，并令 \(\mu_{r_N}\) 与
\[
D_N:=D(r_N)=\sum_{p\in\mathcal P} r_N(p)\,p
\]
如定义 \ref{def:xi-leyang-prime-root-measures} 所示。
假设 \(D_N\to\infty\)，且对每个固定素数 \(p\)，序列 \((r_N(p))_{N}\) 有界。
则：
\begin{enumerate}
  \item 零点经验测度满足 \(\mu_{r_N}\Rightarrow m_{S^1}\)（单位圆上的 Haar 概率测度）；
  \item 与此同时，若记 \(\mathcal A_{r_N}(\theta)=\log|\mathcal L(r_N)(e^{i\theta})|\)，则对每个素数 \(p\) 恒有
  \[
  \frac1\pi\int_{-\pi}^{\pi}\mathcal A_{r_N}(\theta)\cos(p\theta)\,\dd\theta=r_N(p),
  \]
  因而 \((\mathcal A_{r_N})_N\) 不可能在 \(L^2([-\pi,\pi])\) 中收敛到某个与 \(r_N\) 无关的“普适极限”（除非 \(r_N(p)\) 在每个素数处逐点稳定）。
\end{enumerate}
\end{theorem}
\begin{proof}
对任意固定素数 \(q\)，由有界性得 \(r_N(q)\,q=O_q(1)\)，而 \(D_N\to\infty\)，故
\(
\frac{r_N(q)\,q}{D_N}\to 0
\)。
由定理 \ref{thm:xi-leyang-prime-register-haar-criterion} 即得 \(\mu_{r_N}\Rightarrow m_{S^1}\)，从而得到 (1)。
\par
(2) 由定理 \ref{thm:xi-leyang-logamplitude-prime-frequency-inversion} 对每个 \(N\) 与每个素数 \(p\) 的反演恒等式直接给出。
若 \(\mathcal A_{r_N}\to \mathcal A\) 于 \(L^2\)，则其每个余弦系数必收敛，因而 \(r_N(p)\) 必对每个 \(p\) 收敛；若要极限不依赖于序列 \((r_N)\)，则只能在逐素数意义下稳定。证毕。
\end{proof}

\begin{conjecture}[外置诱导素数账本的质量不凝聚与普适差分阶]\label{conj:xi-leyang-prime-register-universal-nonconcentration}
在本文折叠/投影外置范式中，由算法轨迹自然诱导的素数寄存器序列 \((r_m)\) 满足定理 \ref{thm:xi-leyang-prime-register-haar-criterion} 的质量不凝聚条件，从而其 Lee--Yang--Gödel 多项式零点谱 \(\mu_{r_m}\) 弱收敛到 \(m_{S^1}\)。
并且在适当的规范化与均匀性假设下，弧段差分满足 \(\sup_A|\mu_{r_m}(A)-m_{S^1}(A)|=O(D(r_m)^{-1/2})\)，且该指数阶在该范式内具有普适性。
\end{conjecture}

\paragraph{Zeckendorf--Gödel 素数码集合的密度、Lee--Yang 实根结构与谱行列式重整}
设 \((p_k)_{k\ge 1}\) 为递增素数序列，并定义
\[
\mathrm{ZG}
:=\Bigl\{
n\in\NN:\ \forall p,\ v_p(n)\in\{0,1\},\ \forall k\ge 1,\ v_{p_k}(n)\,v_{p_{k+1}}(n)=0
\Bigr\}.
\]
即 \(n\) 为平方因子自由，且其素因子索引集合不含相邻整数。

\begin{theorem}[Gödel--Zeckendorf 素数约束 Dirichlet 级数的对数重整化与半临界解析性]\label{thm:xi-zg-dirichlet-log-renormalization}
\leanverified{paper\_xi\_zg\_dirichlet\_log\_renormalization}
令
\[
\Omega:=\Bigl\{\varepsilon=(\varepsilon_k)_{k\ge 1}:\ \varepsilon_k\in\{0,1\},\ \varepsilon_k\varepsilon_{k+1}=0,\ \varepsilon\ \text{仅有限支撑}\Bigr\}.
\]
对 \(\Re(s)>1\) 定义 Dirichlet 级数
\[
F(s):=\sum_{n\in \mathrm{ZG}}n^{-s}
=\sum_{\varepsilon\in\Omega}\ \prod_{k\ge 1}p_k^{-s\varepsilon_k}.
\]
并对 \(N\ge 0\) 定义截断
\[
F_N(s):=\sum_{\substack{\varepsilon\in\{0,1\}^N\\ \varepsilon_k\varepsilon_{k+1}=0}}\ \prod_{k=1}^{N}p_k^{-s\varepsilon_k},
\qquad
F_0(s)=1,\ \ F_1(s)=1+2^{-s}.
\]
令 \(r_N(s):=F_N(s)/F_{N-1}(s)\)（\(N\ge 1\)），并令
\[
\Delta_N(s):=\log r_N(s)-p_N^{-s}\qquad(N\ge 2),
\]
其中对数取为与实轴上 \(s>1\) 的正值一致的解析分支。
则：
\begin{enumerate}
  \item 对所有 \(\Re(s)>0\) 成立递推恒等式
  \[
  F_N(s)=F_{N-1}(s)+p_N^{-s}F_{N-2}(s),\qquad
  r_N(s)=1+\frac{p_N^{-s}}{r_{N-1}(s)}.
  \]
  \item 级数
  \[
  G(s):=\sum_{N\ge 2}\Delta_N(s)
  \]
  在每个半平面 \(\Re(s)\ge \sigma_0>\tfrac12\) 上绝对且一致收敛，从而在 \(\Re(s)>\tfrac12\) 上给出全纯函数。
  \item 令素数 \(\zeta\) 函数（prime zeta）
  \(
  P(s):=\sum_{k\ge 1}p_k^{-s}
  \)
  （在 \(\Re(s)>1\) 绝对收敛）。
  则对 \(\Re(s)>1\) 有严格分解
  \[
  \log F(s)=P(s)+G(s),
  \qquad\text{从而}\qquad
  F(s)=\exp\!\bigl(P(s)+G(s)\bigr).
  \]
\end{enumerate}
\end{theorem}

\begin{proof}
第一条由末位 \(\varepsilon_N\in\{0,1\}\) 分拆：\(\varepsilon_N=0\) 贡献为 \(F_{N-1}\)，而 \(\varepsilon_N=1\) 强制 \(\varepsilon_{N-1}=0\) 并贡献 \(p_N^{-s}F_{N-2}\)。
对 \(r_N=F_N/F_{N-1}\) 取比即得递推。
\par
对第二条，令 \(\sigma=\Re(s)\ge \sigma_0>\tfrac12\)。
由递推 \(r_N=1+p_N^{-s}/r_{N-1}\) 与 \(|p_N^{-s}|\le p_N^{-\sigma}\) 可在每个紧子域上固定 \(r_{N-1}(s)\) 远离 \(0\) 的一致下界，
从而存在常数 \(C(\sigma_0)\) 使
\[
|\Delta_N(s)|
\le
C(\sigma_0)\Bigl(p_N^{-2\sigma}+(p_Np_{N-1})^{-\sigma}\Bigr)
\qquad(\sigma\ge\sigma_0).
\]
由于
\(
\sum_N p_N^{-2\sigma}
\)
与
\(
\sum_N(p_Np_{N-1})^{-\sigma}
\)
在 \(\sigma>\tfrac12\) 收敛，故 \(\sum_{N\ge 2}\Delta_N(s)\) 在 \(\Re(s)\ge\sigma_0\) 上绝对一致收敛，从而 \(G\) 在 \(\Re(s)>\tfrac12\) 全纯。
\par
第三条由
\(
\log F_N(s)=\sum_{k=1}^N\log r_k(s)
\)
与 \(\Delta_k=\log r_k-p_k^{-s}\) 得
\[
\log F_N(s)=\sum_{k=1}^{N}p_k^{-s}+\sum_{k=2}^{N}\Delta_k(s).
\]
令 \(N\to\infty\) 并在 \(\Re(s)>1\) 用绝对收敛交换极限，得到 \(\log F(s)=P(s)+G(s)\)。证毕。
\end{proof}

\begin{theorem}[平方自由 Euler 因子分离与 \texorpdfstring{$\zeta$}{zeta}-因子化]\label{thm:xi-zg-zeta-factorization}
\leanverified{paper\_xi\_zg\_zeta\_factorization}
沿用定理 \ref{thm:xi-zg-dirichlet-log-renormalization} 的记号，并令
\[
B_N(s):=\prod_{k=1}^{N}\bigl(1+p_k^{-s}\bigr),\qquad
H_N(s):=\frac{F_N(s)}{B_N(s)}.
\]
则存在唯一函数 \(H\) 满足：
\begin{enumerate}
  \item \(H_N(s)\to H(s)\) 在每个紧集 \(K\subset\{s:\Re(s)>\tfrac12\}\) 上一致收敛；
  \item \(H(s)\) 在半平面 \(\Re(s)>\tfrac12\) 上全纯且处处非零；
  \item 在 \(\Re(s)>1\) 上成立严格因子化
  \[
  F(s)=H(s)\frac{\zeta(s)}{\zeta(2s)}.
  \]
\end{enumerate}
从而 \(F\) 在 \(\Re(s)>\tfrac12\) 上亚纯，且在该半平面内唯一可能的极点位于 \(s=1\)，并且为一阶极点。
\end{theorem}

\begin{proof}
由递推可写 \(F_N(s)=\prod_{k=1}^{N}r_k(s)\)，从而
\[
\log H_N(s)
=\sum_{k=1}^{N}\Bigl(\log r_k(s)-\log\bigl(1+p_k^{-s}\bigr)\Bigr).
\]
对 \(k\ge 2\) 使用 \(\Delta_k(s)=\log r_k(s)-p_k^{-s}\)，并将上式改写为
\[
\log H_N(s)
=\Bigl(\log r_1(s)-\log(1+p_1^{-s})\Bigr)
\sum_{k=2}^{N}\Delta_k(s)
\sum_{k=2}^{N}\Bigl(p_k^{-s}-\log(1+p_k^{-s})\Bigr).
\]
第一项恒为 \(0\)。
由定理 \ref{thm:xi-zg-dirichlet-log-renormalization}，\(\sum_{k\ge 2}\Delta_k(s)\) 在 \(\Re(s)>\tfrac12\) 的每个紧集上一致收敛。
另一方面，当 \(\Re(s)\ge \sigma_0>\tfrac12\) 时有
\(
\bigl|p_k^{-s}-\log(1+p_k^{-s})\bigr|
\ll p_k^{-2\sigma_0}
\),
从而 \(\sum_{k\ge 2}\bigl(p_k^{-s}-\log(1+p_k^{-s})\bigr)\) 在同一紧集上一致绝对收敛。
故 \(\log H_N\) 在 \(\Re(s)>\tfrac12\) 的紧集上一致收敛到某个全纯函数 \(\log H\)，并且 \(H=\exp(\log H)\) 全纯且不取零值。

对 \(\Re(s)>1\)，由平方自由 Dirichlet 级数的标准恒等式
\[
\prod_{p}\bigl(1+p^{-s}\bigr)
=\sum_{\substack{m\ge 1\\ m\ \text{平方因子自由}}}\frac1{m^s}
=\frac{\zeta(s)}{\zeta(2s)}
\]
（参见 \cite{Apostol1976,Titchmarsh1986RiemannZeta}），得 \(B_N(s)\to \zeta(s)/\zeta(2s)\)。
再由 \(F_N(s)=H_N(s)B_N(s)\) 并令 \(N\to\infty\)（在 \(\Re(s)>1\) 可由绝对收敛交换极限），得到 \(F(s)=H(s)\zeta(s)/\zeta(2s)\)。
最后，\(\zeta(2s)\) 在 \(\Re(s)>\tfrac12\) 上解析且无零点，而 \(\zeta(s)\) 在该半平面内唯一奇性为 \(s=1\) 的一阶极点，因而 \(F\) 在 \(\Re(s)>\tfrac12\) 上仅可能在 \(s=1\) 处有一阶极点。证毕。
\end{proof}

在平方自由 Euler 因子分离之外，\(\mathrm{ZG}\) 的 Dirichlet 级数还可进一步写成一条递归重整化的乘法链，其局部因子由相邻素数禁配约束逐步修正。

\begin{theorem}[Zeckendorf--Gödel Dirichlet 级数的 continued-Euler 型分解]\label{thm:xi-zg-continued-euler-factorization}
\leanverified{paper\_xi\_zg\_continued\_euler\_factorization}
沿用定理 \ref{thm:xi-zg-dirichlet-log-renormalization} 的记号。对实数 \(s>1\)，定义递归因子
\[
R_1(s):=p_1^{-s},
\qquad
R_n(s):=\frac{p_n^{-s}}{1+R_{n-1}(s)}
\qquad(n\ge 2).
\]
则对每个 \(n\ge 1\) 都有
\[
\boxed{\;
0<R_n(s)<p_n^{-s}.\;}
\]
并且 \(\mathrm{ZG}\) 的 Dirichlet 级数
\[
F(s)=\sum_{n\in \mathrm{ZG}}n^{-s}
\]
具有精确乘法分解
\[
\boxed{\;
F(s)=\prod_{n=1}^{\infty}\bigl(1+R_n(s)\bigr).\;}
\]
特别地，
\[
\boxed{\;
1<F(s)<\prod_{n=1}^{\infty}\bigl(1+p_n^{-s}\bigr)=\frac{\zeta(s)}{\zeta(2s)}
\qquad(s>1).\;}
\]
此外，\(F(s)\) 在区间 \((1,\infty)\) 上严格递减，并满足
\[
\boxed{\;
\lim_{s\to+\infty}F(s)=1.\;}
\]
\end{theorem}
\begin{proof}
由定理 \ref{thm:xi-zg-dirichlet-log-renormalization}，截断和
\[
F_N(s):=\sum_{\substack{\varepsilon\in\{0,1\}^N\\ \varepsilon_k\varepsilon_{k+1}=0}}
\prod_{k=1}^{N}p_k^{-s\varepsilon_k}
\]
满足递推
\[
F_N(s)=F_{N-1}(s)+p_N^{-s}F_{N-2}(s),
\qquad
F_0(s)=1,\quad F_1(s)=1+p_1^{-s}.
\]
再记
\[
r_N(s):=\frac{F_N(s)}{F_{N-1}(s)}
\qquad(N\ge 1).
\]
则同一定理给出
\[
r_1(s)=1+p_1^{-s},
\qquad
r_N(s)=1+\frac{p_N^{-s}}{r_{N-1}(s)}
\qquad(N\ge 2).
\]
定义
\[
R_N(s):=r_N(s)-1.
\]
于是
\[
R_1(s)=p_1^{-s},
\qquad
R_N(s)=\frac{p_N^{-s}}{1+R_{N-1}(s)}
\qquad(N\ge 2),
\]
即得所示递推。由 \(R_1(s)>0\) 出发归纳可知 \(R_N(s)>0\)；再因 \(1+R_{N-1}(s)>1\)，便有
\[
R_N(s)<p_N^{-s}.
\]

另一方面，
\[
F_N(s)=F_{N-1}(s)\,r_N(s)=F_{N-1}(s)\bigl(1+R_N(s)\bigr),
\]
故迭代得到
\[
F_N(s)=\prod_{n=1}^{N}\bigl(1+R_n(s)\bigr).
\]
由于 \(s>1\) 时 \(F_N(s)\to F(s)\)，令 \(N\to\infty\) 即得
\[
F(s)=\prod_{n=1}^{\infty}\bigl(1+R_n(s)\bigr).
\]

由 \(0<R_n(s)<p_n^{-s}\)，逐项比较可得
\[
1+R_n(s)<1+p_n^{-s}.
\]
因此
\[
1<F(s)=\prod_{n\ge 1}\bigl(1+R_n(s)\bigr)
<
\prod_{n\ge 1}\bigl(1+p_n^{-s}\bigr).
\]
最后一项等于 \(\zeta(s)/\zeta(2s)\) 是定理 \ref{thm:xi-zg-zeta-factorization} 中已经使用的平方自由 Euler 乘积恒等式。

再证单调性。若 \(1<s_1<s_2\)，则对每个 \(n\in\mathrm{ZG}\) 有
\[
n^{-s_1}\ge n^{-s_2},
\]
且在 \(n=2\) 处严格不等，于是
\[
F(s_1)-F(s_2)
=
\sum_{n\in\mathrm{ZG}}\bigl(n^{-s_1}-n^{-s_2}\bigr)>0.
\]
故 \(F\) 在 \((1,\infty)\) 上严格递减。

最后，由上界
\[
F(s)<\frac{\zeta(s)}{\zeta(2s)}
\]
且
\[
\frac{\zeta(s)}{\zeta(2s)}\longrightarrow 1
\qquad(s\to+\infty),
\]
并结合显然的下界 \(F(s)>1\)，由夹逼定理得到
\[
\lim_{s\to+\infty}F(s)=1.
\]
证毕。
\end{proof}


\begin{corollary}[素数路径硬核常数与 \texorpdfstring{$s=1$}{s=1} 留数]\label{cor:xi-zg-hardcore-constant-residue}
在定理 \ref{thm:xi-zg-zeta-factorization} 的记号下，记
\[
\mathfrak C_{\mathrm{HC}}:=H(1)=\lim_{N\to\infty}\frac{F_N(1)}{\prod_{k=1}^{N}\bigl(1+\frac1{p_k}\bigr)}\in(0,1).
\]
则 \(F(s)\) 在 \(s=1\) 的留数满足
\[
\operatorname*{Res}_{s=1}F(s)=\frac{\mathfrak C_{\mathrm{HC}}}{\zeta(2)}=\frac{6}{\pi^2}\,\mathfrak C_{\mathrm{HC}}.
\]
并且该留数与 \(\mathrm{ZG}\) 的 Dirichlet 密度一致；结合定理 \ref{con:xi-zg-density-closed-form-tail-bound} 的自然密度存在性，得到
\[
\delta_{\mathrm{ZG}}=\frac{\mathfrak C_{\mathrm{HC}}}{\zeta(2)}.
\]
数值核验可得到
\[
\mathfrak C_{\mathrm{HC}}=0.8461958103\ldots,
\qquad
\delta_{\mathrm{ZG}}=0.5144253666\ldots
\]
\leanverified{paper\_xi\_zg\_hardcore\_constant\_residue}
\end{corollary}

\begin{proof}
由定理 \ref{thm:xi-zg-zeta-factorization}，在 \(s=1\) 的邻域内
\(
F(s)=H(s)\zeta(s)/\zeta(2s)
\)
且 \(H\) 在 \(s=1\) 解析。
因此
\[
\operatorname*{Res}_{s=1}F(s)
=H(1)\operatorname*{Res}_{s=1}\frac{\zeta(s)}{\zeta(2s)}
=\frac{H(1)}{\zeta(2)}.
\]
留数等于 \(\mathrm{ZG}\) 的 Dirichlet 密度是推论 \ref{cor:xi-zg-dirichlet-residue} 的内容，而该推论亦指出留数与自然密度一致，从而得 \(\delta_{\mathrm{ZG}}=\mathfrak C_{\mathrm{HC}}/\zeta(2)\)。
最后，由于对实数 \(s>1\) 有 \(F_N(s)\le B_N(s)\) 且约束非平凡，故 \(0<H(s)<1\) 并推出 \(\mathfrak C_{\mathrm{HC}}\in(0,1)\)。证毕。
\end{proof}

\begin{proposition}[\texorpdfstring{$s=1$}{s=1} 的 Laurent 展开与 Euler--Kronecker 型常数]\label{prop:xi-zg-hardcore-laurent}
记
\[
\kappa_{\mathrm{HC}}:=\left.\frac{d}{ds}\log H(s)\right|_{s=1}.
\]
则 \(\kappa_{\mathrm{HC}}\) 存在，并且 \(F\) 在 \(s=1\) 的 Laurent 展开为
\[
F(s)=\frac{\mathfrak C_{\mathrm{HC}}}{\zeta(2)}\cdot\frac{1}{s-1}
\;+\;
\frac{\mathfrak C_{\mathrm{HC}}}{\zeta(2)}
\left(\gamma-\frac{2\zeta'(2)}{\zeta(2)}+\kappa_{\mathrm{HC}}\right)
 + O(s-1),
\qquad s\to 1,
\]
其中 \(\gamma\) 为 Euler--Mascheroni 常数。
数值核验表明 \(\kappa_{\mathrm{HC}}=0.40404\ldots\)。
\end{proposition}

\begin{proof}
由定理 \ref{thm:xi-zg-zeta-factorization} 的构造，\(\log H_N(s)\) 在 \(\Re(s)>\tfrac12\) 的紧集上一致收敛为 \(\log H(s)\)，并且其尾项由 \(\sum_{k\ge 2}\Delta_k(s)\) 与 \(\sum_{k\ge 2}(p_k^{-s}-\log(1+p_k^{-s}))\) 的一致收敛控制。
同样地，对 \(s\) 位于 \(s=1\) 的小邻域且满足 \(\Re(s)>\tfrac12\)，上述两级数可逐项求导，从而 \(\kappa_{\mathrm{HC}}\) 存在。

另一方面，
\(
\zeta(s)=\frac1{s-1}+\gamma+O(s-1)
\)
且
\[
\frac{1}{\zeta(2s)}
\;=\;
\frac{1}{\zeta(2)}
\left(1-\frac{2\zeta'(2)}{\zeta(2)}(s-1)+O((s-1)^2)\right).
\]
再由
\(
H(s)=H(1)\bigl(1+\kappa_{\mathrm{HC}}(s-1)+O((s-1)^2)\bigr)
\)
三者相乘即得所述展开式。证毕。
\end{proof}

% Auto-split to keep each TeX source < 600 lines.

\begin{theorem}[Zeckendorf--Gödel 素数码集合计数函数的幂节省误差]\label{thm:xi-zg-counting-power-saving-error}
记
\[
A(x):=\#\{n\le x:\ n\in\mathrm{ZG}\}
=\sum_{n\le x}\mathbf 1_{\mathrm{ZG}}(n).
\]
则对任意 \(\varepsilon>0\)，当 \(x\to\infty\) 有
\[
A(x)=\delta_{\mathrm{ZG}}\,x+O_\varepsilon\!\bigl(x^{\frac12+\varepsilon}\bigr).
\]
\end{theorem}

\begin{proof}
令 \(F(s)=\sum_{n\in\mathrm{ZG}}n^{-s}\)。
由定理 \ref{thm:xi-zg-zeta-factorization}，\(F\) 在 \(\Re(s)>\tfrac12\) 上亚纯且在 \(s=1\) 处有一阶极点；
其留数为 \(\delta_{\mathrm{ZG}}\)（推论 \ref{cor:xi-zg-hardcore-constant-residue}）。
\par
对非负系数 Dirichlet 级数应用 Perron 反演（参见 \cite{Apostol1976,Titchmarsh1986RiemannZeta}）可表示 \(A(x)\) 为沿竖线 \(\Re(s)=c>1\) 的积分。
将积分线移至 \(\Re(s)=\tfrac12+\varepsilon\) 时穿越 \(s=1\) 的留数贡献为 \(\delta_{\mathrm{ZG}}x\)。
由于 \(\zeta(2s)\) 在 \(\Re(s)\ge\tfrac12+\varepsilon\) 上解析且远离 \(0\)，并且 \(H(s)\) 在该半平面解析且由其对数展开的绝对收敛给出多项式增长界（定理 \ref{thm:xi-zg-dirichlet-log-renormalization}--\ref{thm:xi-zg-zeta-factorization}），结合 \(\zeta(s)\) 的经典竖线界即可控制移线后的积分余项，从而得到
\(
A(x)=\delta_{\mathrm{ZG}}x+O_\varepsilon(x^{\frac12+\varepsilon})
\)。
证毕。
\end{proof}

\begin{corollary}[Zeckendorf--Gödel 素数码集合的倒数调和渐近]\label{cor:xi-zg-reciprocal-harmonic-asymptotic}
在定理 \ref{thm:xi-zg-counting-power-saving-error} 的记号下，对任意 \(0<\varepsilon<\tfrac12\)，存在常数 \(C_{\mathrm{ZG}}\in\RR\) 使
\[
\sum_{\substack{n\le X\\ n\in\mathrm{ZG}}}\frac{1}{n}
=\delta_{\mathrm{ZG}}\log X+C_{\mathrm{ZG}}+O_\varepsilon\!\bigl(X^{-\frac12+\varepsilon}\bigr)
\qquad(X\to\infty).
\]
特别地，
\[
\sum_{\substack{n\le X\\ n\in\mathrm{ZG}}}\frac{1}{n}
\sim \delta_{\mathrm{ZG}}\log X.
\]
\end{corollary}

\begin{proof}
记误差项
\[
E(x):=A(x)-\delta_{\mathrm{ZG}}x.
\]
由定理 \ref{thm:xi-zg-counting-power-saving-error}，对任意 \(0<\varepsilon<\tfrac12\) 有
\[
E(x)=O_\varepsilon\!\bigl(x^{\frac12+\varepsilon}\bigr).
\]
对计数函数 \(A\) 应用 Abel 求和，得到
\[
\sum_{\substack{n\le X\\ n\in\mathrm{ZG}}}\frac{1}{n}
=\frac{A(X)}{X}+\int_{1}^{X}\frac{A(t)}{t^{2}}\,\dd t.
\]
代入 \(A(t)=\delta_{\mathrm{ZG}}t+E(t)\) 可得
\[
\sum_{\substack{n\le X\\ n\in\mathrm{ZG}}}\frac{1}{n}
=\delta_{\mathrm{ZG}}\log X+\delta_{\mathrm{ZG}}
+\frac{E(X)}{X}
+\int_{1}^{X}\frac{E(t)}{t^{2}}\,\dd t.
\]
由于
\[
\frac{E(X)}{X}=O_\varepsilon\!\bigl(X^{-\frac12+\varepsilon}\bigr),
\qquad
\frac{E(t)}{t^{2}}=O_\varepsilon\!\bigl(t^{-\frac32+\varepsilon}\bigr),
\]
且指数 \(-\tfrac32+\varepsilon<-1\)，故积分
\[
\int_{1}^{\infty}\frac{E(t)}{t^{2}}\,\dd t
\]
绝对收敛。定义
\[
C_{\mathrm{ZG}}
:=\delta_{\mathrm{ZG}}+\int_{1}^{\infty}\frac{E(t)}{t^{2}}\,\dd t.
\]
于是
\[
\int_{1}^{X}\frac{E(t)}{t^{2}}\,\dd t
=C_{\mathrm{ZG}}-\delta_{\mathrm{ZG}}+O_\varepsilon\!\bigl(X^{-\frac12+\varepsilon}\bigr),
\]
从而得到所述渐近展开。特别地，主项为 \(\delta_{\mathrm{ZG}}\log X\)，故结论成立。
\end{proof}



\begin{theorem}[Zeckendorf--Gödel 素数码集合的自然密度、闭式可计算性与尾项误差界]\label{con:xi-zg-density-closed-form-tail-bound}
对每个 \(N\ge 1\)，定义截断集合
\[
\mathrm{ZG}^{(N)}
:=\Bigl\{
n\in\NN:\ \forall 1\le k\le N,\ v_{p_k}(n)\in\{0,1\},\
\forall 1\le k\le N-1,\ v_{p_k}(n)\,v_{p_{k+1}}(n)=0
\Bigr\}.
\]
记
\[
\delta_N:=\mathrm{dens}\bigl(\mathrm{ZG}^{(N)}\bigr).
\]
则自然密度
\[
\delta_{\mathrm{ZG}}:=\mathrm{dens}(\mathrm{ZG})
\]
存在，且
\[
\delta_{\mathrm{ZG}}=\lim_{N\to\infty}\delta_N\in(0,1).
\]
进一步，\(\delta_N\) 具有显式闭式
\[
\delta_N=\Bigl(\prod_{k=1}^N\Bigl(1-\frac1{p_k}\Bigr)\Bigr)I_N,
\qquad
I_N:=\sum_{\substack{S\subseteq\{1,\dots,N\}\\ S\ \mathrm{无相邻}}}\ \prod_{k\in S}\frac1{p_k},
\]
并满足尾项误差上界
\[
0\le \delta_N-\delta_{\mathrm{ZG}}
\le
\sum_{k>N}\frac1{p_k^2}
\;+\;
\sum_{k\ge N}\frac1{p_kp_{k+1}}.
\]
特别地，\(\delta_{\mathrm{ZG}}\) 可通过生成素数并迭代计算 \(\delta_N\) 在任意精度下有效逼近。
\end{theorem}
\begin{proof}
对固定 \(N\)，成员资格仅依赖于模
\(
M_N:=\prod_{k=1}^N p_k^2
\)
的剩余类，故 \(\delta_N\) 存在。

记
\(
a_k:=1-\frac1{p_k},\ b_k:=a_k\frac1{p_k}
\)。
在前 \(N\) 个素数坐标上，“允许指数为 \(1\)”的索引集合恰为路径图 \(\{1,\dots,N\}\) 的独立集 \(S\)；
对应权重为
\(
\prod_{k\in S}b_k\prod_{k\notin S}a_k
\)。
求和得
\[
\delta_N
=\sum_{S\ \mathrm{无相邻}}\ \prod_{k\in S}b_k\prod_{k\notin S}a_k
=\Bigl(\prod_{k=1}^N a_k\Bigr)\sum_{S\ \mathrm{无相邻}}\prod_{k\in S}\frac1{p_k},
\]
即所示闭式。

又 \(\mathrm{ZG}=\bigcap_{N\ge 1}\mathrm{ZG}^{(N)}\)，故 \(\delta_N\) 单调下降且
\[
\overline{\mathrm{dens}}(\mathrm{ZG})\le \delta_N\quad(\forall N).
\]
若 \(n\in \mathrm{ZG}^{(N)}\setminus \mathrm{ZG}\)，则必发生尾部违约事件之一：
\[
\exists\,k>N,\ p_k^2\mid n
\qquad\text{或}\qquad
\exists\,k\ge N,\ p_kp_{k+1}\mid n.
\]
故
\[
\mathrm{ZG}^{(N)}\setminus \mathrm{ZG}
\subseteq
\Bigl(\bigcup_{k>N}\{p_k^2\mid n\}\Bigr)
\cup
\Bigl(\bigcup_{k\ge N}\{p_kp_{k+1}\mid n\}\Bigr).
\]
由并集上界与
\(
\mathrm{dens}\{m\mid n\}=1/m
\)
得
\[
\mathrm{dens}\bigl(\mathrm{ZG}^{(N)}\setminus \mathrm{ZG}\bigr)
\le
\sum_{k>N}\frac1{p_k^2}
+
\sum_{k\ge N}\frac1{p_kp_{k+1}},
\]
即误差估计。令 \(N\to\infty\) 得上下密度同限，故自然密度存在且等于 \(\lim\delta_N\)。

最后，
\[
\delta_{\mathrm{ZG}}
\ge 1-\sum_{k\ge 1}\frac1{p_k^2}-\sum_{k\ge 1}\frac1{p_kp_{k+1}}>0,
\]
因为两级数收敛且其和严格小于 \(1\)。证毕。
\end{proof}

\begin{corollary}[Mertens 重整常数与素数路径独立集分区函数的对数增长律]\label{cor:xi-zg-mertens-log-growth}
沿用定理 \ref{con:xi-zg-density-closed-form-tail-bound} 的记号。
则
\[
\delta_{\mathrm{ZG}}
=\lim_{N\to\infty}
\Bigl(\prod_{k=1}^N\Bigl(1-\frac1{p_k}\Bigr)\Bigr)I_N,
\]
并且
\[
I_N\sim \delta_{\mathrm{ZG}}\,e^{\gamma}\,\log p_N
\qquad (N\to\infty),
\]
其中 \(\gamma\) 为 Euler--Mascheroni 常数。
\end{corollary}
\begin{proof}
由定理 \ref{con:xi-zg-density-closed-form-tail-bound}，
\[
I_N=\frac{\delta_N}{\prod_{k\le N}(1-\frac1{p_k})}.
\]
再用 Mertens 乘积渐近
\(
\prod_{p\le y}(1-\frac1p)\sim e^{-\gamma}/\log y
\)
并令 \(y=p_N\)、\(\delta_N\to\delta_{\mathrm{ZG}}\) 即得。证毕。
\end{proof}

\begin{theorem}[素数权重 Fibonacci 路径的 Lee--Yang 实根性与 Poisson--binomial 分解]\label{con:xi-zg-weighted-path-leyang-poisson-binomial}
定义
\[
P_N(t):=\sum_{\substack{S\subseteq\{1,\dots,N\}\\ S\ \mathrm{无相邻}}}
t^{|S|}\prod_{k\in S}\frac1{p_k},
\]
则
\[
P_0(t)=1,\qquad
P_1(t)=1+\frac{t}{p_1},\qquad
P_N(t)=P_{N-1}(t)+\frac{t}{p_N}P_{N-2}(t)\quad(N\ge 2).
\]
并且成立：
\begin{enumerate}
  \item \(P_N\) 的全部零点均为互异负实根；
  \item \(P_N\) 与 \(P_{N-1}\) 的零点严格交错；
  \item 设 \(m_N:=\deg P_N=\lfloor(N+1)/2\rfloor\)，则存在 \(\alpha_{N,j}>0\) 使
  \[
  P_N(t)=\prod_{j=1}^{m_N}\Bigl(1+\frac{t}{\alpha_{N,j}}\Bigr).
  \]
\end{enumerate}
对固定 \(t>0\)，令
\[
\pi_{N,j}(t):=\frac{t}{t+\alpha_{N,j}}\in(0,1).
\]
若按 Gibbs 权重
\(
\propto t^{|S|}\prod_{k\in S}p_k^{-1}
\)
在独立集上取样，则随机变量 \(|S|\) 满足分布分解
\[
|S|\overset{d}{=}\sum_{j=1}^{m_N}B_{N,j},
\qquad
B_{N,j}\sim\mathrm{Bernoulli}\!\bigl(\pi_{N,j}(t)\bigr),
\]
其中 \((B_{N,j})\) 相互独立。
\end{theorem}
\begin{proof}
递推直接由索引 \(N\) 是否入独立集的分拆得到。

令 \(\mathcal T_{N+1}\) 为长度 \(N+1\) 的路径树，赋边权
\(
w_k:=p_k^{-1/2}
\)
（\(1\le k\le N\)），并记 \(A_{N+1}\) 为其加权邻接矩阵。
匹配多项式与独立集母函数的标准对应给出
\[
\chi_{N+1}(x):=\det(xI-A_{N+1})
=x^{N+1}P_N\!\Bigl(-\frac1{x^2}\Bigr).
\]
由于 \(A_{N+1}\) 为实对称 Jacobi 矩阵（副对角元严格正），其特征值实且互异；
故 \(P_N\) 的零点为 \(-1/\lambda_j^2\)，均为互异负实数。
交错性由主子矩阵谱交错（Cauchy interlacing）传递到 \(P_N,P_{N-1}\) 的零点。

第三条由实负根分解直接成立。
设
\[
G_{N,t}(s):=\frac{P_N(ts)}{P_N(t)}
=\prod_{j=1}^{m_N}
\frac{1+\frac{ts}{\alpha_{N,j}}}{1+\frac{t}{\alpha_{N,j}}}
=\prod_{j=1}^{m_N}\bigl(1-\pi_{N,j}(t)+\pi_{N,j}(t)s\bigr),
\]
右侧即独立 Bernoulli 和的概率母函数，故得到分布恒等式。证毕。
\end{proof}

\paragraph{素数权硬核归一化常数、Mertens--Lee--Yang 渐近与系数不等式链}
沿用定理 \ref{con:xi-zg-weighted-path-leyang-poisson-binomial} 的素数权路径独立集分区函数 \(P_N(t)\) 与其递推。
\par
\begin{theorem}[硬核归一化常数的存在性与显式尾项控制]\label{thm:xi-prime-hardcore-normalization-constant}
\leanverified{paper\_xi\_prime\_hardcore\_normalization\_constant}
对 \(t>0\) 与 \(N\ge 1\) 定义归一化比值
\begin{equation}\label{eq:xi-prime-hardcore-RN-def}
R_N(t)
:=
\frac{P_N(t)}{\prod_{k=1}^{N}\left(1+\frac{t}{p_k}\right)}\in(0,1].
\end{equation}
则：
\begin{enumerate}
  \item \((R_N(t))_{N\ge 1}\) 严格单调递减；
  \item 极限
  \begin{equation}\label{eq:xi-prime-hardcore-deltaHC-def}
  \boxed{\ \delta_{\mathrm{HC}}(t):=\lim_{N\to\infty}R_N(t)\ }
  \end{equation}
  存在且满足 \(0<\delta_{\mathrm{HC}}(t)<1\)；
  \item 对所有 \(N\ge 2\) 有尾项界
  \begin{equation}\label{eq:xi-prime-hardcore-RN-tail}
  \boxed{\;
  0\ \le\ R_N(t)-\delta_{\mathrm{HC}}(t)\ \le\ \sum_{k\ge N}\frac{t^2}{p_kp_{k+1}}\; }.
  \end{equation}
\end{enumerate}
\end{theorem}
\begin{proof}
记 \(a_k:=t/p_k\) 与 \(D_N(t):=\prod_{k=1}^{N}(1+a_k)\)。
由递推 \(P_N=P_{N-1}+a_NP_{N-2}\) 得
\[
R_N=\frac{R_{N-1}}{1+a_N}+\frac{a_N}{(1+a_{N-1})(1+a_N)}R_{N-2}.
\]
对 \(N\ge 3\) 进一步消元得到差分恒等式
\[
R_{N-1}-R_N
=\frac{a_{N-1}a_N}{(1+a_{N-2})(1+a_{N-1})(1+a_N)}\,R_{N-3}\ >\ 0,
\]
从而 \((R_N)\) 严格递减。
又由于 \(P_N\) 计数的独立集严格包含于所有子集，故 \(0<P_N\le D_N\)，从而 \(0<R_N\le 1\)。
\par
由上式立得上界 \(0<R_{N-1}-R_N\le a_{N-1}a_N=t^2/(p_{N-1}p_N)\)。
由于 \(\sum_{N\ge 2}p_{N-1}^{-1}p_N^{-1}\) 收敛，故 \(\sum_{N\ge 2}(R_{N-1}-R_N)\) 收敛，进而 \(R_N\) 收敛并满足尾项界 \eqref{eq:xi-prime-hardcore-RN-tail}。
\par
最后证明极限严格为正。
由 \(R_{N-2}\ge R_{N-1}\) 与上一差分恒等式得
\[
R_N
\ge R_{N-1}\left(1-\frac{a_{N-1}a_N}{(1+a_{N-1})(1+a_N)}\right),
\]
对 \(N\) 连乘并用 \(\sum_{N\ge 2}a_{N-1}a_N<\infty\) 得无穷乘积下界严格为正，从而 \(\delta_{\mathrm{HC}}(t)>0\)。证毕。
\end{proof}

\begin{corollary}[\texorpdfstring{$t=1$}{t=1} 与 \texorpdfstring{$\mathfrak C_{\mathrm{HC}}$}{C_HC} 的一致性]\label{cor:xi-prime-hardcore-deltaHC-matches-C_HC}
\leanverified{paper\_xi\_prime\_hardcore\_deltahc\_matches\_c\_hc}
在推论 \ref{cor:xi-zg-hardcore-constant-residue} 的记号下，有
\[
\delta_{\mathrm{HC}}(1)=\mathfrak C_{\mathrm{HC}}.
\]
从而
\(
\delta_{\mathrm{ZG}}=\delta_{\mathrm{HC}}(1)/\zeta(2)
\)
与推论 \ref{cor:xi-zg-hardcore-constant-residue} 的闭式一致。
\end{corollary}
\begin{proof}
由 \eqref{eq:xi-prime-hardcore-RN-def} 在 \(t=1\) 的定义，\(\delta_{\mathrm{HC}}(1)=\lim_{N\to\infty}P_N(1)/\prod_{k\le N}(1+1/p_k)\)。
而推论 \ref{cor:xi-zg-hardcore-constant-residue} 中 \(\mathfrak C_{\mathrm{HC}}\) 正是同一极限。证毕。
\end{proof}

\begin{theorem}[Mertens--Lee--Yang 渐近与两常数分解]\label{thm:xi-prime-hardcore-mertens-leyang-asymptotic}
令 \(\delta_{\mathrm{HC}}(t)\) 如定理 \ref{thm:xi-prime-hardcore-normalization-constant} 定义，并定义收敛欧拉乘积
\begin{equation}\label{eq:xi-prime-hardcore-Ct-euler}
C(t):=\prod_{p}\Bigl(1-\frac{1}{p}\Bigr)^{t}\Bigl(1+\frac{t}{p}\Bigr)\in(0,\infty).
\end{equation}
则当 \(N\to\infty\) 时有
\begin{equation}\label{eq:xi-prime-hardcore-asymptotic}
\boxed{\;
P_N(t)\ \sim\ \delta_{\mathrm{HC}}(t)\,C(t)\,e^{\gamma t}\,(\log p_N)^{t}\;}
\end{equation}
其中 \(\gamma\) 为 Euler--Mascheroni 常数。
\end{theorem}
\begin{proof}
由 \eqref{eq:xi-prime-hardcore-RN-def} 与 \eqref{eq:xi-prime-hardcore-deltaHC-def}，
\(
P_N(t)=R_N(t)\prod_{k\le N}(1+t/p_k)\sim \delta_{\mathrm{HC}}(t)\prod_{k\le N}(1+t/p_k)
\)。
将有限乘积分解为
\[
\prod_{k\le N}\Bigl(1+\frac{t}{p_k}\Bigr)
=
\Bigl(\prod_{k\le N}\Bigl(1-\frac{1}{p_k}\Bigr)^{-t}\Bigr)
\cdot
\Bigl(\prod_{k\le N}\Bigl(1-\frac{1}{p_k}\Bigr)^t\Bigl(1+\frac{t}{p_k}\Bigr)\Bigr).
\]
第二个括号内因子满足
\(
\bigl(1-\frac1p\bigr)^t(1+\frac{t}{p})
=1-\frac{t(t+1)}{2p^2}+O(p^{-3})
\)，
其对数为 \(O(p^{-2})\)，故乘积收敛并趋于 \(C(t)\)。
第一个括号由 Mertens 乘积渐近
\(
\prod_{p\le y}(1-\frac1p)\sim e^{-\gamma}/\log y
\)
（取 \(y=p_N\)）给出
\(
\prod_{k\le N}(1-\frac1{p_k})^{-t}\sim e^{\gamma t}(\log p_N)^t
\)。
合并即得 \eqref{eq:xi-prime-hardcore-asymptotic}。证毕。
\end{proof}

\begin{theorem}[素数权硬核 Gibbs 场的规模涨落与中心极限定理]\label{thm:xi-prime-hardcore-gibbs-clt}
\leanverified{paper\_xi\_prime\_hardcore\_gibbs\_clt}
固定 \(t>0\)。
在 \(\{1,\dots,N\}\) 的无相邻集合上按权重
\(
\propto t^{|S|}\prod_{k\in S}p_k^{-1}
\)
抽样，并记 \(K_N:=|S|\)。
则存在仅依赖于 \(t\) 的常数 \(\mu(t),\nu(t)\in\RR\)，使得
\[
\EE[K_N]=t\log\log p_N+\mu(t)+o(1),
\qquad
\mathrm{Var}(K_N)=t\log\log p_N+\nu(t)+o(1).
\]
并且令 \(V_N:=\mathrm{Var}(K_N)\)，则 \(V_N\to\infty\) 且存在绝对常数 \(C>0\) 使
\[
\sup_{x\in\RR}\left|
\PP\!\left(\frac{K_N-\EE[K_N]}{\sqrt{V_N}}\le x\right)-\Phi(x)
\right|
\le \frac{C}{\sqrt{V_N}}
=O\!\left(\frac{1}{\sqrt{\log\log p_N}}\right),
\]
从而 \(\frac{K_N-\EE[K_N]}{\sqrt{V_N}}\Rightarrow\mathcal N(0,1)\)。
\end{theorem}
\begin{proof}
由定理 \ref{thm:xi-prime-hardcore-mertens-leyang-asymptotic}，
\[
\log P_N(t)=t\log\log p_N+\log\delta_{\mathrm{HC}}(t)+\log C(t)+\gamma t+o(1).
\]
对该指数族分配，有标准恒等式
\(
\EE[K_N]=t\,\partial_t\log P_N(t)
\)
与
\(
\mathrm{Var}(K_N)=t\,\partial_t\EE[K_N]
\)，
从而得到所示两条渐近；其中常数 \(\mu(t),\nu(t)\) 由 \(\log\delta_{\mathrm{HC}}(t)+\log C(t)+\gamma t\) 的一、二阶导数给出。
\par
另一方面，由定理 \ref{con:xi-zg-weighted-path-leyang-poisson-binomial} 的 Poisson--binomial 分解，
\(
K_N\overset{d}{=}\sum_{j=1}^{m_N}B_{N,j}
\)
为独立 Bernoulli 和，且 \(\sum_j\mathrm{Var}(B_{N,j})=V_N\to\infty\)。
对 Poisson--binomial 和应用 Berry--Esseen 型界（以第三绝对中心矩上界为代价）得到误差 \(O(V_N^{-1/2})\)，从而得到中心极限定理。证毕。
\end{proof}

\begin{corollary}[单位权素数硬核 Gibbs 场的中心极限定律]\label{cor:xi-prime-hardcore-gibbs-clt-t1-explicit}
\leanverified{paper\_xi\_prime\_hardcore\_gibbs\_clt\_t1\_explicit}
在定理 \ref{thm:xi-prime-hardcore-gibbs-clt} 中取 \(t=1\)。
记
\[
I_N
:=
\sum_{\substack{S\subseteq\{1,\dots,N\}\\ S\ \mathrm{无相邻}}}\prod_{i\in S}\frac1{p_i},
\qquad
\PP_N(S):=\frac{1}{I_N}\prod_{i\in S}\frac1{p_i},
\]
并令 \(\omega_N:=|S|\)。
则
\[
\frac{\omega_N-\mu_N}{\sigma_N}\Longrightarrow \mathcal N(0,1),
\]
其中
\[
\mu_N=\sum_{i=1}^{N}\frac{1}{p_i+1}+O(1),
\qquad
\sigma_N^2=\sum_{i=1}^{N}\frac{p_i}{(p_i+1)^2}+O(1),
\]
并且
\[
\sigma_N^2\sim \sum_{i\le N}\frac1{p_i}\sim \log\log p_N.
\]
\end{corollary}
\begin{proof}
定理 \ref{thm:xi-prime-hardcore-gibbs-clt} 在 \(t=1\) 时给出
\[
\EE[\omega_N]=\log\log p_N+O(1),
\qquad
\mathrm{Var}(\omega_N)=\log\log p_N+O(1),
\]
并且
\(
(\omega_N-\EE[\omega_N])/\sqrt{\mathrm{Var}(\omega_N)}
\Rightarrow
\mathcal N(0,1)
\)。

另一方面，
\[
\frac1{p_i}-\frac{1}{p_i+1}=\frac{1}{p_i(p_i+1)}=O(p_i^{-2}),
\]
以及
\[
\frac1{p_i}-\frac{p_i}{(p_i+1)^2}
=
\frac{2p_i+1}{p_i(p_i+1)^2}
=
O(p_i^{-2}).
\]
由于 \(\sum_i p_i^{-2}<\infty\)，故
\[
\sum_{i=1}^{N}\frac{1}{p_i+1}
=
\sum_{i=1}^{N}\frac1{p_i}+O(1),
\qquad
\sum_{i=1}^{N}\frac{p_i}{(p_i+1)^2}
=
\sum_{i=1}^{N}\frac1{p_i}+O(1).
\]
再由素数倒数和的 Mertens 渐近，
\[
\sum_{i=1}^{N}\frac1{p_i}=\log\log p_N+O(1),
\]
即得所示 \(\mu_N,\sigma_N^2\) 的表达式与发散结论。
最后，由 \(\sigma_N^2=\mathrm{Var}(\omega_N)+O(1)\) 且 \(\sigma_N^2\to\infty\)，两种标准化只相差 \(1+o(1)\) 因子，从而极限正态性保持。证毕。
\end{proof}

\begin{proposition}[素数寄存器与 Lee--Yang 谱的 Newton--Maclaurin 对应及不等式链]\label{prop:xi-prime-hardcore-newton-maclaurin}
\leanverified{paper\_xi\_prime\_hardcore\_newton\_maclaurin}
设 \(m_N=\deg P_N=\lfloor(N+1)/2\rfloor\)，并写
\[
P_N(t)=\sum_{k=0}^{m_N}a_{N,k}\,t^k.
\]
则：
\begin{enumerate}
  \item 系数具有素数寄存器刻画
  \[
  a_{N,k}
  =\sum_{\substack{S\subseteq\{1,\dots,N\}\\ S\ \mathrm{无相邻},\ |S|=k}}
  \ \prod_{i\in S}\frac1{p_i}.
  \]
  \item 若 \(P_N(t)=\prod_{j=1}^{m_N}(1+t/\alpha_{N,j})\)（\(\alpha_{N,j}>0\)；定理 \ref{con:xi-zg-weighted-path-leyang-poisson-binomial}），则
  \[
  a_{N,k}=e_k\!\left(\frac{1}{\alpha_{N,1}},\dots,\frac{1}{\alpha_{N,m_N}}\right),
  \]
  其中 \(e_k\) 为第 \(k\) 个初等对称多项式；并且
  \[
  \sum_{j=1}^{m_N}\frac{1}{\alpha_{N,j}}=\sum_{i=1}^{N}\frac{1}{p_i}.
  \]
  \item 序列
  \[
  \left(\frac{a_{N,k}}{\binom{m_N}{k}}\right)^{1/k}\qquad(k=1,\dots,m_N)
  \]
  随 \(k\) 单调不增；并且 \(a_{N,k}\) 满足严格 Newton 不等式 \(a_{N,k}^2>a_{N,k-1}a_{N,k+1}\)（\(1\le k\le m_N-1\)）。
\end{enumerate}
\end{proposition}
\begin{proof}
前两条由 \(P_N\) 的独立集展开与 Lee--Yang 实根分解展开逐项比较得到。
第三条由正数列的 Maclaurin 不等式与实负互异零点多项式的严格 Newton 不等式推出。证毕。
\end{proof}


\begin{theorem}[Lee--Yang 单位圆提升、Joukowsky 椭圆输运与 \texorpdfstring{$\delta_{\mathrm{ZG}}$}{deltaZG} 的谱行列式表达]\label{con:xi-zg-unit-circle-joukowsky-spectral-determinant}
在定理 \ref{con:xi-zg-weighted-path-leyang-poisson-binomial} 的记号下：
\begin{enumerate}
  \item 设 \(\chi_{N+1}(x)=\det(xI-A_{N+1})\)，\(\Phi_{N+1}(x):=\chi_{N+1}(x/\sqrt2)\)。
  由于 \(\|A_{N+1}\|_{2}\le 2\max_k w_k=\sqrt2\)，\(\Phi_{N+1}\) 全部根落在 \([-2,2]\)。
  因此
  \[
  \mathcal L_{N+1}(z)
  :=z^{N+1}\Phi_{N+1}(z+z^{-1})
  =z^{N+1}\chi_{N+1}\!\Bigl(\frac{z+z^{-1}}{\sqrt2}\Bigr)
  \]
  的全部零点位于单位圆 \(|z|=1\)。
  \item 对任意 \(r>1\)，令 \(J_r(z):=rz+r^{-1}z^{-1}\)。
  若 \(z_j\) 为 \(\mathcal L_{N+1}\) 的零点，则 \(w_j:=J_r(z_j)\) 落在椭圆 \(J_r(S^1)\) 上。
  \item 取 \(x=i\) 于恒等式
  \(
  \chi_{N+1}(x)=x^{N+1}P_N(-x^{-2})
  \)
  得
  \[
  P_N(1)=|\chi_{N+1}(i)|=|\det(iI-A_{N+1})|.
  \]
  因而
  \[
  \delta_{\mathrm{ZG}}
  =\lim_{N\to\infty}
  \Bigl(\prod_{k=1}^N\Bigl(1-\frac1{p_k}\Bigr)\Bigr)\,|\det(iI-A_{N+1})|.
  \]
  又
  \[
  |\det(iI-A_{N+1})|^2
  =\det\!\bigl((iI-A_{N+1})(-iI-A_{N+1})\bigr)
  =\det(I+A_{N+1}^2),
  \]
  故
  \[
  P_N(1)=\det(I+A_{N+1}^2)^{1/2},
  \quad
  \delta_{\mathrm{ZG}}
  =\lim_{N\to\infty}
  \Bigl(\prod_{k=1}^N\Bigl(1-\frac1{p_k}\Bigr)\Bigr)\det(I+A_{N+1}^2)^{1/2}.
  \]
\end{enumerate}
\end{theorem}
\begin{proof}
第一条由谱半径界和参数化 \(x=z+z^{-1}=2\cos\theta\) 直接得到。
第二条是 \(J_r\) 对单位圆的标准像性质。
第三条由 \(x=i\) 代入与行列式恒等式逐步化简得出。证毕。
\end{proof}

\begin{corollary}[\texorpdfstring{$s=1$}{s=1} 处单极点与 Dirichlet 留数常数]\label{cor:xi-zg-dirichlet-residue}
令 \(P(s)=\sum_{k\ge 1}p_k^{-s}\) 为 prime zeta，并令 \(G(s)\) 为定理 \ref{thm:xi-zg-dirichlet-log-renormalization} 的正则因子。
采用 prime zeta 在 \(s\to 1^+\) 的经典展开
\[
P(s)=\log\frac1{s-1}+B_{\mathsf P}+o(1),
\]
其中 \(B_{\mathsf P}\in\RR\) 为常数（可由 Mertens 定理推导；参见 \cite{Apostol1976,Titchmarsh1986RiemannZeta}）。
则 \(F(s)=\sum_{n\in\mathrm{ZG}}n^{-s}\) 在 \(s=1\) 处具有单极点，且其留数为
\[
\lim_{s\to1^+}(s-1)F(s)=\exp\!\bigl(B_{\mathsf P}+G(1)\bigr)\in(0,\infty).
\]
并且该留数等于 \(\mathrm{ZG}\) 的 Dirichlet 密度；结合定理 \ref{con:xi-zg-density-closed-form-tail-bound} 的自然密度存在性，两者一致。
\end{corollary}

\paragraph{Abel 边界律与前缀递推证书}

\begin{theorem}[ZG 的 Dirichlet 级数在 \texorpdfstring{$s=1$}{s=1} 的 Abel 型边界律、调和主项与收敛临界线]\label{thm:xi-zg-abel-residue-log-density}
沿用本段记号，记
\[
A_{\mathrm{ZG}}(x):=\#\{n\le x:\ n\in\mathrm{ZG}\},
\qquad
D_{\mathrm{ZG}}(s):=\sum_{n\in\mathrm{ZG}}n^{-s}
\qquad (\Re(s)>1),
\]
以及
\[
\delta_{\mathrm{ZG}}:=\lim_{x\to\infty}\frac{A_{\mathrm{ZG}}(x)}{x}.
\]
则 \(\delta_{\mathrm{ZG}}\in(0,1)\)，并且成立：
\begin{enumerate}
  \item Abel 型边界律
  \[
  \lim_{\sigma\downarrow 1}(\sigma-1)D_{\mathrm{ZG}}(\sigma)=\delta_{\mathrm{ZG}}.
  \]
  \item 调和和主项
  \[
  \sum_{\substack{n\le X\\ n\in\mathrm{ZG}}}\frac{1}{n}
  =\delta_{\mathrm{ZG}}\log X+o(\log X)
  \qquad (X\to\infty).
  \]
  \item 绝对收敛临界线满足
  \[
  \sigma_c(D_{\mathrm{ZG}})=1.
  \]
\end{enumerate}
\end{theorem}
\begin{proof}
由定理 \ref{con:xi-zg-density-closed-form-tail-bound}，自然密度存在且
\[
\delta_{\mathrm{ZG}}\in(0,1).
\]
记
\[
A_{\mathrm{ZG}}(x)=\delta_{\mathrm{ZG}}\,x+r(x),
\qquad
r(x)=o(x)
\qquad (x\to\infty).
\]

对任意 \(\sigma>1\)，Abel--Stieltjes 分部积分给出
\[
D_{\mathrm{ZG}}(\sigma)
=
\sigma\int_1^\infty A_{\mathrm{ZG}}(x)x^{-\sigma-1}\,\dd x.
\]
代入 \(A_{\mathrm{ZG}}(x)=\delta_{\mathrm{ZG}}x+r(x)\) 后得到
\[
D_{\mathrm{ZG}}(\sigma)
=
\frac{\sigma\,\delta_{\mathrm{ZG}}}{\sigma-1}
+
\sigma\int_1^\infty r(x)x^{-\sigma-1}\,\dd x.
\]
因此只需证明
\[
\lim_{\sigma\downarrow1}
(\sigma-1)\int_1^\infty r(x)x^{-\sigma-1}\,\dd x=0.
\]
任取 \(\varepsilon>0\)。
由 \(r(x)=o(x)\)，可取 \(X_0\ge 1\) 使得 \(|r(x)|\le \varepsilon x\) 对所有 \(x\ge X_0\) 成立。
于是
\[
\begin{aligned}
(\sigma-1)\left|\sigma\int_1^\infty r(x)x^{-\sigma-1}\,\dd x\right|
&\le
(\sigma-1)\sigma\int_1^{X_0}|r(x)|x^{-\sigma-1}\,\dd x \\
&\qquad
+(\sigma-1)\sigma\int_{X_0}^{\infty}\varepsilon x\cdot x^{-\sigma-1}\,\dd x \\
&\le
(\sigma-1)C_{X_0}
+\sigma\varepsilon X_0^{1-\sigma},
\end{aligned}
\]
其中 \(C_{X_0}<\infty\) 与 \(\sigma\) 无关。
令 \(\sigma\downarrow1\)，上式上极限不超过 \(\varepsilon\)；再令 \(\varepsilon\downarrow0\)，即得
\[
\lim_{\sigma\downarrow 1}(\sigma-1)D_{\mathrm{ZG}}(\sigma)=\delta_{\mathrm{ZG}}.
\]

再看第二条。
对和式作 Abel 分部积分，
\[
\sum_{\substack{n\le X\\ n\in\mathrm{ZG}}}\frac1n
=
\frac{A_{\mathrm{ZG}}(X)}{X}+\int_1^X \frac{A_{\mathrm{ZG}}(t)}{t^2}\,\dd t.
\]
代入 \(A_{\mathrm{ZG}}(t)=\delta_{\mathrm{ZG}}t+r(t)\) 得
\[
\sum_{\substack{n\le X\\ n\in\mathrm{ZG}}}\frac1n
=
\delta_{\mathrm{ZG}}
+\frac{r(X)}{X}
+\delta_{\mathrm{ZG}}\log X
+\int_1^X \frac{r(t)}{t^2}\,\dd t.
\]
由 \(r(t)=o(t)\)，对任意 \(\varepsilon>0\) 仍可取 \(T_0\) 使 \(|r(t)|\le \varepsilon t\) 对 \(t\ge T_0\) 成立，从而
\[
\left|\int_1^X \frac{r(t)}{t^2}\,\dd t\right|
\le C_{T_0}+\varepsilon\log X,
\]
其中 \(C_{T_0}\) 与 \(X\) 无关。
因此
\[
\sum_{\substack{n\le X\\ n\in\mathrm{ZG}}}\frac1n
=
\delta_{\mathrm{ZG}}\log X+o(\log X).
\]
最后，由 \(\mathrm{ZG}\subset \NN\) 可知对任意 \(\sigma>1\) 都有
\[
0\le D_{\mathrm{ZG}}(\sigma)\le \zeta(\sigma)<\infty,
\]
故 \(D_{\mathrm{ZG}}\) 在半平面 \(\Re(s)>1\) 上绝对收敛，从而 \(\sigma_c(D_{\mathrm{ZG}})\le 1\)。另一方面，第二条等价于
\[
\sum_{\substack{n\le X\\ n\in\mathrm{ZG}}}\frac1n
=
\delta_{\mathrm{ZG}}\log X+o(\log X)\to\infty,
\]
从而级数 \(D_{\mathrm{ZG}}(1)=\sum_{n\in\mathrm{ZG}}n^{-1}\) 发散。于是其收敛 abscissa 恰为 \(1\)。
\end{proof}

\begin{corollary}[ZG 的 Dirichlet 级数的绝对收敛临界线恰为 \texorpdfstring{$\Re(s)=1$}{Re(s)=1}]\label{cor:xi-zg-absolute-convergence-critical-line}
在定理 \ref{thm:xi-zg-abel-residue-log-density} 的设定下，\(D_{\mathrm{ZG}}(s)\) 在 \(\Re(s)>1\) 上绝对收敛，在 \(s=1\) 发散，并满足
\[
D_{\mathrm{ZG}}(\sigma)\sim \frac{\delta_{\mathrm{ZG}}}{\sigma-1}
\qquad (\sigma\downarrow 1).
\]
因而 \(\Re(s)=1\) 正是其绝对收敛临界线。
\end{corollary}
\begin{proof}
由定理 \ref{thm:xi-zg-abel-residue-log-density} 的第一条，
\[
(\sigma-1)D_{\mathrm{ZG}}(\sigma)\longrightarrow \delta_{\mathrm{ZG}}>0
\qquad (\sigma\downarrow 1),
\]
故
\[
D_{\mathrm{ZG}}(\sigma)
=
\frac{\delta_{\mathrm{ZG}}}{\sigma-1}
+o\!\bigl((\sigma-1)^{-1}\bigr)
\qquad (\sigma\downarrow 1).
\]
特别地，\(D_{\mathrm{ZG}}(\sigma)\to+\infty\)。再结合同一定理的第三条，便得到所述绝对收敛临界线结论。证毕。
\end{proof}

\begin{proposition}[Zeckendorf--Gödel 密度的归一化前缀递推、单调收敛与尾误差证书]\label{prop:xi-zg-prefix-recursion-monotone-tail-certificate}
沿用定理 \ref{con:xi-zg-density-closed-form-tail-bound} 与
\ref{con:xi-zg-weighted-path-leyang-poisson-binomial} 的记号。
定义数列
\[
(a_0,b_0)=(1,0),
\]
并对 \(N\ge 1\) 递推地令
\[
a_N=(a_{N-1}+b_{N-1})\frac{p_N}{p_N+1},
\qquad
b_N=\frac{a_{N-1}}{p_N+1}.
\]
再记
\[
u_N:=a_N+b_N,
\qquad
\widetilde\delta_N:=\frac{u_N}{\zeta(2)}.
\]
则成立：
\begin{enumerate}
  \item 对每个 \(N\ge 0\)，有
  \[
  u_N=
  \frac{P_N(1)}{\prod_{k=1}^{N}\left(1+\frac1{p_k}\right)},
  \]
  其中约定空积等于 \(1\)。
  因而
  \[
  \delta_{\mathrm{ZG}}
  =
  \frac{1}{\zeta(2)}\lim_{N\to\infty}u_N.
  \]
  \item \(u_0=u_1=1\)，并且数列 \((\widetilde\delta_N)_{N\ge 1}\) 严格单调下降，满足
  \[
  \widetilde\delta_N\downarrow \delta_{\mathrm{ZG}}.
  \]
  \item 对每个 \(N\ge 1\)，有显式尾误差界
  \[
  0\le \widetilde\delta_N-\delta_{\mathrm{ZG}}
  \le
  \widetilde\delta_N
  \sum_{j=N}^{\infty}\frac{1}{(p_j+1)(p_{j+1}+1)}.
  \]
\end{enumerate}
\end{proposition}
\begin{proof}
对 \(N\ge 1\)，定义两类部分配分函数
\[
A_N:=
\sum_{\substack{S\subseteq\{1,\dots,N\}\\ S\ \mathrm{无相邻}\\ N\notin S}}
\prod_{k\in S}\frac1{p_k},
\qquad
B_N:=
\sum_{\substack{S\subseteq\{1,\dots,N\}\\ S\ \mathrm{无相邻}\\ N\in S}}
\prod_{k\in S}\frac1{p_k}.
\]
则
\[
P_N(1)=A_N+B_N.
\]
并且由末位是否被选中，直接得到
\[
A_N=P_{N-1}(1),
\qquad
B_N=\frac{1}{p_N}A_{N-1}.
\]
把它们按
\[
a_N:=\frac{A_N}{\prod_{k=1}^{N}\left(1+\frac1{p_k}\right)},
\qquad
b_N:=\frac{B_N}{\prod_{k=1}^{N}\left(1+\frac1{p_k}\right)}
\]
归一化，便得到所述递推关系
\[
a_N=(a_{N-1}+b_{N-1})\frac{p_N}{p_N+1},
\qquad
b_N=\frac{a_{N-1}}{p_N+1}.
\]
于是
\[
u_N=a_N+b_N
=
\frac{P_N(1)}{\prod_{k=1}^{N}\left(1+\frac1{p_k}\right)}.
\]
这就证明了第一条中的第一式。

另一方面，由定理 \ref{con:xi-zg-density-closed-form-tail-bound}，
\[
\delta_N
=
\left(\prod_{k=1}^{N}\left(1-\frac1{p_k}\right)\right)P_N(1).
\]
将上一式代入并整理得
\[
\delta_N
=
\left(\prod_{k=1}^{N}\left(1-\frac1{p_k^2}\right)\right)u_N.
\]
由于 \(\delta_N\to \delta_{\mathrm{ZG}}\) 且
\[
\prod_{k=1}^{N}\left(1-\frac1{p_k^2}\right)\longrightarrow \frac1{\zeta(2)},
\]
只要 \(u_N\) 收敛，便有
\[
\delta_{\mathrm{ZG}}=\frac{1}{\zeta(2)}\lim_{N\to\infty}u_N.
\]

现证单调性。
由递推直接计算
\[
u_N
=
u_{N-1}-\frac{b_{N-1}}{p_N+1}.
\]
因此 \(u_0=u_1=1\)。
而对 \(N\ge 2\)，由 \(a_{N-2}>0\) 可得
\[
b_{N-1}=\frac{a_{N-2}}{p_{N-1}+1}>0,
\]
故
\[
u_N<u_{N-1}\qquad (N\ge 2).
\]
于是 \((u_N)_{N\ge 1}\) 严格递减且下有界于 \(0\)，故极限
\(
u_\infty:=\lim_{N\to\infty}u_N
\)
存在。
由前述关系，便得到
\[
\widetilde\delta_N=\frac{u_N}{\zeta(2)}\downarrow \frac{u_\infty}{\zeta(2)}=\delta_{\mathrm{ZG}}.
\]

最后证明尾误差界。
由
\[
u_N-u_{N+1}=\frac{b_N}{p_{N+1}+1}
\]
以及
\[
b_N=\frac{a_{N-1}}{p_N+1},
\qquad
u_N=a_{N-1}+\frac{p_N}{p_N+1}b_{N-1}\ge a_{N-1},
\]
可得
\[
u_N-u_{N+1}
\le
\frac{u_N}{(p_N+1)(p_{N+1}+1)}.
\]
而由于 \((u_N)\) 递减，对任意 \(j\ge N\) 亦有
\[
u_j-u_{j+1}
\le
\frac{u_j}{(p_j+1)(p_{j+1}+1)}
\le
\frac{u_N}{(p_j+1)(p_{j+1}+1)}.
\]
故
\[
u_N-u_\infty
=
\sum_{j=N}^{\infty}(u_j-u_{j+1})
\le
u_N\sum_{j=N}^{\infty}\frac{1}{(p_j+1)(p_{j+1}+1)}.
\]
两边乘以 \(\zeta(2)^{-1}\) 即得
\[
0\le \widetilde\delta_N-\delta_{\mathrm{ZG}}
\le
\widetilde\delta_N\sum_{j=N}^{\infty}\frac{1}{(p_j+1)(p_{j+1}+1)}.
\]
证毕。
\end{proof}


\begin{proposition}[临界重整化常数的可计算误差界]\label{prop:xi-zg-holographic-constant-tail-bound}
\leanverified{paper\_xi\_zg\_holographic\_constant\_tail\_bound}
取 \(s=1\)，并令 \(\Delta_N(1)\) 与 \(G(1)=\sum_{N\ge 2}\Delta_N(1)\) 如定理 \ref{thm:xi-zg-dirichlet-log-renormalization}。
定义截断
\[
G^{(N)}(1):=\sum_{k=2}^{N}\Delta_k(1).
\]
则存在绝对常数 \(C>0\) 使对所有 \(N\ge 3\) 有
\[
\bigl|G(1)-G^{(N)}(1)\bigr|
\le
C\sum_{k>N}\Bigl(\frac1{p_k^2}+\frac1{p_kp_{k-1}}\Bigr),
\]
从而正则因子 \(\exp(G(1))\) 可被一条显式可求和尾项在任意精度下有效逼近。
\end{proposition}

\begin{proof}
由定理 \ref{thm:xi-zg-dirichlet-log-renormalization} 的一致估计，在 \(s=1\) 时存在常数 \(C\) 使
\(
|\Delta_k(1)|
\le C\bigl(p_k^{-2}+(p_kp_{k-1})^{-1}\bigr)
\)。
对尾和求并界即得。证毕。
\end{proof}

\begin{theorem}[\(\mathrm{ZG}\) 的对数 Mertens 定律]\label{thm:xi-zg-log-mertens}
设
\[
H(x):=\sum_{\substack{n\le x\\ n\in\mathrm{ZG}}}\frac1n.
\]
则极限存在并满足
\[
\lim_{x\to\infty}\frac{H(x)}{\log x}=\delta_{\mathrm{ZG}},
\]
并且常数项可由 \(F(s)=\sum_{n\in\mathrm{ZG}}n^{-s}\) 在 \(s=1\) 的 Laurent 常数项显式识别：存在 \(C_{\mathrm{ZG}}\in\RR\) 使
\[
H(x)=\delta_{\mathrm{ZG}}\log x+C_{\mathrm{ZG}}+o(1)\qquad(x\to\infty).
\]
其中
\[
C_{\mathrm{ZG}}
=\frac{\mathfrak C_{\mathrm{HC}}}{\zeta(2)}
\left(\gamma-\frac{2\zeta'(2)}{\zeta(2)}+\kappa_{\mathrm{HC}}\right)
=\delta_{\mathrm{ZG}}
\left(\gamma-\frac{2\zeta'(2)}{\zeta(2)}+\kappa_{\mathrm{HC}}\right).
\]
\end{theorem}

\begin{proof}
由推论 \ref{cor:xi-zg-dirichlet-residue} 或推论 \ref{cor:xi-zg-hardcore-constant-residue}，\(F(s)=\sum_{n\in\mathrm{ZG}}n^{-s}\) 在 \(s=1\) 处的留数为 \(\delta_{\mathrm{ZG}}\)。
对非负系数 Dirichlet 级数应用标准 Tauberian 定理可得
\(
H(x)=\delta_{\mathrm{ZG}}\log x+C_{\mathrm{ZG}}+o(1)
\)，其中 \(C_{\mathrm{ZG}}\) 等于 \(F\) 在 \(s=1\) 的 Laurent 常数项。
而命题 \ref{prop:xi-zg-hardcore-laurent} 给出了该 Laurent 常数项的显式表达式，从而得到所示恒等式。证毕。
\end{proof}

\begin{conjecture}[正则因子的全半平面延拓与 Lee--Yang 型零点边界]\label{conj:xi-zg-leyang-godel-transport-boundary}
令 \(\mathcal R(s):=\exp(G(s))\)（定理 \ref{thm:xi-zg-dirichlet-log-renormalization}），其在 \(\Re(s)>\tfrac12\) 全纯且无零。
猜想 \(\mathcal R\) 可延拓为 \(\Re(s)>0\) 上的全纯无零函数；从而 \(F(s)=\exp(P(s))\mathcal R(s)\) 在 \(\Re(s)>0\) 的全部奇性结构由 prime zeta \(P(s)\) 的奇性传输唯一决定。
进一步地，对充分大的截断 \(N\)，将 \(F_N(s)\) 视为有限体积分配函数，则其归一化零点云在 \(N\to\infty\) 的极限中形成一条由 Joukowsky 椭圆输运诱导的边界曲线，并以 \(\Re(s)=\tfrac12\) 为临界对称轴，与 \(\mathcal R\) 的零点真空相容。
\end{conjecture}

\begin{conjecture}[弱局域约束下的 Erd\H{o}s--Kac 型高斯极限]\label{conj:xi-zg-erdos-kac-under-local-hardcore}
记
\[
\omega(n):=\#\{p:\ p\mid n\},
\qquad
\mathcal Z(X):=\{n\le X:\ n\in\mathrm{ZG}\}.
\]
猜想存在常数 \(c_0\in\RR\)，使得在 \(\mathcal Z(X)\) 上均匀抽样的 \(n\) 满足
\[
\frac{\omega(n)-\log\log X-c_0}{\sqrt{\log\log X}}
\Longrightarrow
\mathcal N(0,1)
\qquad(X\to\infty).
\]
该猜想表明：相邻素数不可同现的局域硬核约束在中心极限定律尺度上仅导致有限平移，而不改变标准化极限分布。
\end{conjecture}

\begin{remark}[可复现核验脚本]\label{rem:xi-zg-density-leyang-spectral-audit-script}
定理 \ref{con:xi-zg-density-closed-form-tail-bound}--\ref{con:xi-zg-unit-circle-joukowsky-spectral-determinant}
中的闭式恒等式与有限窗数值核验由脚本
\texttt{scripts/exp\_xi\_prime\_register\_zg\_density\_leyang\_spectrum\_audit.py}
给出，对应证书文件为
\texttt{artifacts/export/xi\_prime\_register\_zg\_density\_leyang\_spectrum\_audit.json}。
\par
定理 \ref{thm:xi-zg-dirichlet-log-renormalization} 与命题 \ref{prop:xi-zg-holographic-constant-tail-bound}
在 \(s=1\) 处的递推与尾项模板核验由脚本
\texttt{scripts/exp\_xi\_prime\_register\_zg\_dirichlet\_log\_renormalization\_audit.py}
给出，对应证书文件为
\texttt{artifacts/export/xi\_prime\_register\_zg\_dirichlet\_log\_renormalization\_audit.json}。
\par
定理 \ref{thm:xi-zg-zeta-factorization}--命题 \ref{prop:xi-zg-hardcore-laurent}
中 \(\mathfrak C_{\mathrm{HC}}\) 与 \(\kappa_{\mathrm{HC}}\) 的数值核验由脚本
\texttt{scripts/exp\_xi\_prime\_register\_zg\_hardcore\_zeta\_factorization\_constant\_audit.py}
给出，对应证书文件为
\texttt{artifacts/export/xi\_prime\_register\_zg\_hardcore\_zeta\_factorization\_constant\_audit.json}。
\end{remark}


\begin{theorem}[行列式素因子分解给出的椭圆面积守恒律]\label{thm:xi-terminal-ellipse-area-euler-prime-law}
\leanverified{paper\_xi\_terminal\_ellipse\_area\_euler\_prime\_law}
设 \(A\in\mathrm{GL}_2(\QQ)\)，单位圆盘 \(D:=\{x\in\RR^2:\|x\|\le 1\}\)，椭圆盘 \(E_A:=A(D)\)。
则
\[
\mathrm{Area}(E_A)=\pi\,|\det A|.
\]
并且若
\(
\det A=\pm\prod_p p^{v_p(\det A)}
\)
（仅有限个 \(v_p\neq 0\)），则
\[
\mathrm{Area}(E_A)=\pi\prod_p p^{v_p(\det A)}.
\]
\end{theorem}
\begin{proof}
二维 Lebesgue 测度在可逆线性变换下按雅可比因子 \( |\det A| \) 缩放，故
\(
\mathrm{Area}(A(D))=|\det A|\,\mathrm{Area}(D)=\pi|\det A|
\)。
其余结论由有理数乘法群的素因子分解直接得到。证毕。
\end{proof}

\begin{theorem}[有理椭圆合同类的二元不变量完备性]\label{thm:xi-terminal-rational-ellipse-two-invariant-completeness}
\leanverified{paper\_xi\_terminal\_rational\_ellipse\_two\_invariant\_completeness}
设 \(A\in\mathrm{GL}_2(\QQ)\)，记 \(E_A:=A(S^1)\subset\RR^2\)。
定义
\[
\Delta_1(A):=(\det A)^2\in\QQ_{>0},
\qquad
\Delta_2(A):=\mathrm{tr}(A^\top A)\in\QQ_{>0}.
\]
则 \(E_A\) 的欧氏合同类由 \((\Delta_1(A),\Delta_2(A))\) 完全决定。
若 \(\sigma_1\ge \sigma_2>0\) 为 \(A\) 的奇异值，则
\[
\sigma_1^2,\sigma_2^2\ \text{是}\ t^2-\Delta_2(A)t+\Delta_1(A)=0\ \text{的两根},
\]
因而
\[
\sigma_{1,2}
=\sqrt{\frac{\Delta_2(A)\pm\sqrt{\Delta_2(A)^2-4\Delta_1(A)}}{2}}.
\]
特别地，对任意 \(A,B\in\mathrm{GL}_2(\QQ)\),
\[
E_A\cong E_B\ (\text{欧氏合同})
\Longleftrightarrow
\bigl(\Delta_1(A),\Delta_2(A)\bigr)=\bigl(\Delta_1(B),\Delta_2(B)\bigr).
\]
\end{theorem}
\begin{proof}
极分解 \(A=UP\)（\(U\) 正交，\(P=(A^\top A)^{1/2}\) 对称正定）给出
\(
E_A=U(P(S^1))
\)；
故合同类由 \(P\) 决定。
\(P\) 的特征值为 \(\sigma_1,\sigma_2\)，即椭圆半轴长。
又
\[
\sigma_1^2+\sigma_2^2=\mathrm{tr}(A^\top A)=\Delta_2(A),
\qquad
\sigma_1^2\sigma_2^2=\det(A^\top A)=(\det A)^2=\Delta_1(A),
\]
因此 \((\sigma_1,\sigma_2)\) 由 \((\Delta_1,\Delta_2)\) 唯一恢复，反之亦然。证毕。
\end{proof}

% AUTO-GENERATED by scripts/exp_xi_prime_register_semidirect_two_shadow_audit.py
\begin{table}[H]
\centering
\scriptsize
\setlength{\tabcolsep}{6pt}
\caption{有理矩阵样本的椭圆二元不变量审计：$\Delta_1=(\det A)^2$ 与 $\Delta_2=\mathrm{tr}(A^\top A)$。}
\label{tab:xi_elliptic_two_shadow_invariants_audit}
\begin{tabular}{l l l l r r}
\toprule
id & $\det A$ & $\Delta_1(A)$ & $\Delta_2(A)$ & $\sigma_1$ (num.) & $\sigma_2$ (num.)\\
\midrule
$A_1$ & $209/120$ & $43681/14400$ & $14701/3600$ & 1.762791958 & 0.988016004\\
$B_1=A_1Q_{90}$ & $209/120$ & $43681/14400$ & $14701/3600$ & 1.762791958 & 0.988016004\\
$A_2$ & $163/70$ & $26569/4900$ & $233209/44100$ & 1.974001870 & 1.179619667\\
\bottomrule
\end{tabular}
\end{table}


\begin{theorem}[有限素部格与实部椭圆影的二影刚性]\label{thm:xi-terminal-finite-archimedean-two-shadow-rigidity}
定义 prime-register 纤维
\[
\widehat{\ZZ}:=\varprojlim_m \ZZ/m\ZZ\cong\prod_p \ZZ_p.
\]
对 \(A\in\mathrm{GL}_2(\QQ)\)，记有限素部格影
\[
\Lambda_A:=A\,\widehat{\ZZ}^{\,2}\subset \mathbb A_f^{\,2},
\]
并记实部椭圆影 \(E_A:=A(S^1)\subset\RR^2\)。
若 \(A,B\in\mathrm{GL}_2(\QQ)\) 满足
\[
\Lambda_A=\Lambda_B,
\qquad
E_A=E_B,
\]
则存在 \(Q\in O_2(\ZZ)\) 使 \(A=BQ\)。
若再施加定向保持规范（等价于 \(\det(B^{-1}A)=1\)），则 \(Q\in SO_2(\ZZ)\)，因此仅余 \(4\) 个离散候选；
进一步施加任一破缺该四元对称的规范锚定后，\(Q\) 唯一，从而 \(A\) 由 \((\Lambda_A,E_A)\) 唯一确定。
\end{theorem}
\begin{proof}
令 \(C:=B^{-1}A\)。
由 \(\Lambda_A=\Lambda_B\) 得
\(
C\,\widehat{\ZZ}^{\,2}=\widehat{\ZZ}^{\,2}
\)，故 \(C\in\mathrm{GL}_2(\widehat{\ZZ})\)。

由 \(E_A=E_B\) 得同一椭圆二次型：
\[
x^\top(AA^\top)^{-1}x=1
\Longleftrightarrow
x^\top(BB^\top)^{-1}x=1,
\]
从而 \(AA^\top=BB^\top\)。
左乘 \(B^{-1}\)、右乘 \((B^{-1})^\top\) 得
\[
CC^\top=I,
\]
即 \(C\in O_2(\QQ)\)。

同时 \(C\in\mathrm{GL}_2(\widehat{\ZZ})\) 蕴含其每个有理条目在所有 \(\ZZ_p\) 中整，故条目属于
\(
\QQ\cap\bigcap_p\ZZ_p=\ZZ
\)；
即 \(C\in\mathrm{GL}_2(\ZZ)\)。
综上 \(C\in O_2(\ZZ)\)，即 \(A=BQ\)（\(Q:=C\in O_2(\ZZ)\)）。

当 \(\det Q=1\) 时 \(Q\in SO_2(\ZZ)\)。
二维整正交群中 \(|SO_2(\ZZ)|=4\)，故仅有四个定向保持候选；对任意固定规范锚定（例如长轴正向锚定）即可消去残余离散歧义并得到唯一性。证毕。
\end{proof}
